\documentclass[11pt,a4paper]{article}
\usepackage[T1]{fontenc}
\usepackage[utf8]{inputenc}
\usepackage{amsmath,amsthm}
\usepackage{newpxtext,newpxmath}
\usepackage[a4paper,margin=25mm,headheight=14pt]{geometry}
\usepackage{microtype,mathtools}
\usepackage{booktabs,longtable,array,enumitem}
\usepackage{xcolor,listings,fancyhdr,needspace}
\usepackage{xurl}
\usepackage[colorlinks=true,linkcolor=black,citecolor=black,urlcolor=black,pdfusetitle]{hyperref}
\usepackage{bookmark}
\definecolor{codebg}{gray}{0.965}
\lstset{basicstyle=\small\ttfamily,backgroundcolor=\color{codebg},columns=fullflexible,
 keepspaces=true,breaklines=true,frame=single,framerule=0pt,xleftmargin=5pt,xrightmargin=5pt,
 aboveskip=10pt,belowskip=10pt,showstringspaces=false}
\setlist{nosep,leftmargin=1.5em}
\setlength{\parskip}{4pt}
\setlength{\parindent}{1.2em}
\setlength{\emergencystretch}{2em}
\renewcommand{\arraystretch}{1.16}
\newcolumntype{L}[1]{>{\raggedright\arraybackslash}p{#1}}
\newtheorem{proposition}{Proposition}[section]
\newtheorem{theorem}[proposition]{Theorem}
\newtheorem{lemma}[proposition]{Lemma}
\theoremstyle{definition}
\newtheorem{definition}[proposition]{Definition}
\newtheorem{example}[proposition]{Example}
\theoremstyle{remark}
\newtheorem{remark}[proposition]{Remark}
\newcommand{\code}[1]{\texttt{\detokenize{#1}}}
\newcommand{\List}{\operatorname{List}}
\newcommand{\Option}{\operatorname{Option}}
\newcommand{\foldr}{\operatorname{foldr}}
\newcommand{\rev}{\operatorname{reverse}}
\newcommand{\Nat}{\mathbb N}
\newcommand{\Bool}{\mathbb B}
\newcommand{\concat}{\mathbin{+\!+}}
\newcommand{\ans}{\noindent\textbf{Answer.}\ }
\newcommand{\imp}{\par\noindent\textbf{Implementation consequence.}\ }
\newcommand{\lit}{\par\noindent\textbf{Research evidence.}\ }
\newcommand{\risk}{\par\noindent\textbf{Boundary of the claim.}\ }
\newcommand{\qref}[1]{\hyperref[q:#1]{#1}}
\pagestyle{fancy}
\fancyhf{}
\fancyhead[L]{\small Leant2: answers to the expert questions}
\fancyhead[R]{\small Research analysis}
\fancyfoot[C]{\thepage}
\renewcommand{\headrulewidth}{0.3pt}
\title{\textbf{From Expert Questions to a\linebreak Verifiable Search Architecture}\\[6pt]
\large A research analysis of Leant2's further-improvement proposals}
\author{Research report prepared for the Leant2 project}
\date{21 September 2026}
\hypersetup{pdftitle={From Expert Questions to a Verifiable Search Architecture},
 pdfauthor={Research report prepared for the Leant2 project},
 pdfsubject={Answers to all 31 expert questions in the Leant2 further-improvement proposal},
 pdfkeywords={Lean, Leant2, synthesis, dependent types, carriers, proof search}}
\begin{document}
\maketitle
\thispagestyle{empty}
\begin{center}
\small Repository snapshot: \code{784664ff34e819422510a4d470ab07fe48bb736b}\\
Inspected toolchain: Lean \code{v4.34.0}
\end{center}
\begin{abstract}
This report answers all 31 questions in the nine research areas of Leant2's
\emph{Further Improvements} proposal. The central recommendation is to separate
kernel-certified acceptance, solution-preserving pruning, and reproducible
search and ranking: none of these guarantees implies the other two. Concrete
answers include dependency-separated search over persistent Lean states;
bounded joint synthesis of carriers, algebras, and finishers; residual
observation graphs with rollback-safe invalidation; native motive construction;
proof-producing evaluation; and dictionary-aware library retrieval.

Several premises of the proposal require correction. Reverse is an ordinary
right fold into lists, although an endomorphism carrier can improve its cost.
Finite observations do not determine the minimal carrier of an unknown
function, and partial observation rows cannot simply be counted as states.
Counting open goals is not generally an admissible search heuristic. Smyth
already supports higher-order examples in its implementation, Hoogle+ already
handles type classes, and Lean 4.34 provides \code{cbv} and \code{decide_cbv}
for proof-producing evaluation. The report gives explicit correctness
arguments, bounded completeness statements, implementation interfaces,
regression tests, and an evaluation plan. Included Python experiments check
small semantic claims; no Lean compilation, Leant2 benchmark rerun, or model
accuracy experiment is claimed.
\end{abstract}
\vfill
\noindent\textbf{Reading guide.} Sections~\ref{sec:r1}--\ref{sec:r9} follow the
proposal's R1--R9 numbering. Every question has its own identifier, such as
R2.3. Sections~\ref{sec:architecture} and~\ref{sec:roadmap} combine the answers
into an implementation sequence. Appendix~\ref{app:map} is a complete question
index; Appendix~\ref{app:checks} records the executable checks.
\clearpage
\tableofcontents
\clearpage

\section{Scope, evidence, and the guarantees that matter}\label{sec:scope}

\subsection{What was inspected}
The requested proposal is the directory
\code{docs/proposals/11-further-improvements} in the pinned repository
snapshot~\cite{repo}. Its question-bearing files are
\code{05-search.tex} through \code{12-ranking.tex}; the last contains both R8
and R9. The expert-contact summary in \code{15-experts.tex} groups those
questions rather than introducing another independent list. There are
$4+5+4+4+3+4+3+2+2=31$ questions.

The analysis also inspected the acceptance gate, transactional state, and
engine implementation~\cite{implementation}. This matters because a design
recommendation must fit the actual Lean-native architecture, not a generic
synthesis system. The current engine uses metavariable-state snapshots,
iterative deepening, wall-clock lane deadlines, and a short grace period after
the first accepted candidate. Its ranking uses unused binders, eliminator
counts, and expression size, with printed-form deduplication. The acceptance
path kernel-checks declarations synchronously and audits their transitive
axiom dependencies.

The proposal's performance tables are \emph{reported measurements}, not new
experiments in this report. In particular, its acceptance table is explicitly
anchored to the earlier commit \code{33cec8b}, not to the snapshot inspected
here. Its approximately 45\% residual-evaluation share and approximately 1\%
state-save/restore share come from one contract-heavy query; they do not prove
that large best-first frontiers have negligible memory cost~\cite{proposal-evidence}.
The distinction is important when prioritizing a rewrite.

\subsection{Three independent guarantees}
Let a frozen query be $Q=(\Gamma,T,C,\mathcal A,\mathcal G)$, where $T$ is the
requested type, $C$ is an optional contract, $\mathcal A$ is the permitted axiom
profile, and $\mathcal G$ records the selected search grammar and bounds.
The following are different claims.

\paragraph{Acceptance correctness.}
An accepted result contains a term $p$ and, when required, a proof $\pi$, with
\[
 \Gamma\vdash p:T,\qquad \Gamma\vdash \pi:C(p),
 \qquad \operatorname{axioms}(p,\pi)\subseteq\mathcal A.
\]
The kernel checks the judgments; a request-binding layer must ensure that
these are the \emph{requested} judgments. A kernel cannot detect that a caller
silently substituted a different target or contract.

\paragraph{Pruning correctness.}
For a partial program $p$ with a dependent telescope of holes $\Delta$, a hard
refutation must establish that no permitted completion can satisfy the query:
\begin{equation}\label{eq:refutation}
 \forall\theta:\Delta,\quad B(\theta)\longrightarrow\neg C(p[\theta]).
\end{equation}
Here $B$ records the branch's genuine assumptions, not speculative model
predictions. A final kernel gate does not repair the loss of a correct program
that was wrongly pruned. Conversely, an evaluator may be very useful while
remaining incomplete: it can return \emph{unknown} without endangering
soundness.

\paragraph{Search and publication reproducibility.}
A fixed logical-work budget, grammar, environment, proposal manifest, and
scheduler should determine the same search trace and published candidate
order. A fixed wall-clock deadline generally determines only a
machine-dependent prefix of that trace. Kernel correctness says nothing about
which valid result appears first.

These distinctions should appear in data types and receipts, not merely in
comments. A result may be certified but not optimal; a branch may be unresolved
rather than refuted; and a bounded grammar may be exhausted without proving
that the original Lean type is uninhabited.

\subsection{Seven immediate corrections to the proposal}
The detailed arguments below imply seven changes before implementing the
larger research plan.

First, replace the claim that one unit per open non-leaf goal is an admissible
heuristic with a proved lower bound or zero. Shared unification can discharge
several goals at once. Second, distinguish minimal \emph{semantic state},
minimal \emph{type expression}, and minimal or efficient \emph{program}; these
are different carrier objectives. Third, do not use reverse as evidence that
a list-valued right fold is impossible. Fourth, require an actual upper bound
on reachable carrier states before using a cardinality refutation; counting
values found so far supplies the opposite kind of bound.

Fifth, \code{n + 0} is definitionally equal to \code{n} in Lean; an example
requiring transport should use \code{0 + n} or another genuinely
non-definitional equality~\cite{lean-recursion}. Sixth, observational
coincidence is a useful display relation but is not generally a congruence for
higher-order or recursive synthesis. Seventh, update the related-work picture:
Smyth, Hoogle+, and Lean's current evaluator already answer important parts of
R3, R5, and R7~\cite{smyth,hoogle,lean-cbv}.

\subsection{Decision overview}
\begin{longtable}{@{}L{9mm}L{53mm}L{89mm}@{}}
\toprule
Area & Recommended first implementation & What must not be claimed\\
\midrule\endhead
R1 & Persistent whole-state queue with dependency components & Independent goals from target syntax alone; machine-independent wall-clock results\\
R2 & Bounded joint carrier--algebra--finisher synthesis & A globally minimal intended carrier learned from finite examples\\
R3 & Observation DAG plus proof-producing refutation backend & Differential tests as a proof of conservativity\\
R4 & Native induction adapter and dependency-closed generalization & Complete general motive inference; arbitrary cast erasure\\
R5 & Typed recursive capabilities with checked realization & Automatic definitional equality between two interpreters\\
R6 & Local invariant obligations and watcher-driven proof debt & A complete syntactic recognizer for inductively provable properties\\
R7 & Typed hypergraph retrieval with explicit dictionaries & Sound negative reachability for rules omitted from the abstraction\\
R8 & Frozen publication objective and layered equivalence & Sample equivalence as unrestricted semantic equality\\
R9 & Optional recorded proposals with deterministic replay & Temperature zero as a reproducibility guarantee; unmeasured accuracy\\
\bottomrule
\end{longtable}

\section{R1: search with coupled goals}\label{sec:r1}

\subsection{R1.1: can shared goals be organized without copying everything?}\label{q:R1.1}
\ans Yes: use a persistent \emph{whole-state OR graph}, and factor only certified
independent parts of a state into AND components. This does not eliminate
alternative assignments to genuinely shared variables. It avoids treating
all goals as either globally independent or globally sequential.

A state should contain the exact Lean saved state, the pending goals, delayed
constraints, residual observations, immutable provenance, and the costs needed
by the scheduler. Build a dependency graph whose vertices include expression
metavariables, universe metavariables, and pending constraints. A goal's support
is the transitive closure through its target, local declaration types and
values, assignments, and delayed dependencies. A variable absent from the
printed target can still be reachable through a local's type. A shared
universe constraint can couple apparently unrelated data goals.

For two candidate components, compute or conservatively approximate their
read and write footprints. If either can assign a variable that the other
reads, they belong in the same component. Unknown tactic effects also prevent
separation. Shared \emph{fixed} locals need not connect components; shared
revisable state does. A small separator of shared variables can be assigned
first, after which the remaining components often become independent.

\begin{proposition}[Conditional factorization]\label{prop:factor}
Suppose a state's obligations partition into components with disjoint mutable
sets $W_1,\ldots,W_k$. Every transition for component $i$ reads only immutable
shared data and $W_i$, writes only $W_i$ (including fresh variables allocated
exclusively to that component), and leaves all cross-component
constraints unchanged. Then transitions in different components commute, and
successful component solutions can be combined in any order.
\end{proposition}
\begin{proof}
A transition on $W_i$ cannot change the inputs observed by a transition on
$W_j$ for $i\ne j$, so the two compositions have the same result on each
mutable set and on shared data. Repeatedly exchange adjacent independent
transitions to group a successful interleaving by component. Conversely,
concatenating successful component traces preserves every component's
preconditions. The assumptions ensure no omitted joint obligation remains.
\end{proof}

This theorem is an interface contract, not a claim that arbitrary Lean tactics
already satisfy it. Start with conservative footprints for native synthesis
rules; treat a general tactic as potentially global unless sandboxed.

\lit Aesop provides the relevant practical comparison: its treatment of shared
metavariables includes copying and reapplying affected goals rather than
assuming independent subproofs. Canonical-min provides a more recent precedent
for postponing constraints on blockers and resuming them under assignments in
persistent, backtracking search. Its own calculus is not a replacement for
Lean's authoritative expressions~\cite{aesop,canonical-min}.

\imp Memoize a solved component as a telescope-abstracted, typed term template,
with its environment and assumption fingerprint. Instantiate it with fresh
universe and expression variables, substitute actual locals, and recheck its
type. Never move raw free-variable identifiers between unrelated contexts.
Cache failure only with a bounded-exhaustion justification or a checked
refutation, not because an earlier proof attempt timed out.

\subsection{R1.2: what heuristic is genuinely admissible?}\label{q:R1.2}
\ans The universally safe starting point is $h(s)=0$. Useful nonzero bounds
must be lower bounds on \emph{joint completion cost}, not counts of obligations.

Consider two obligations that become reflexive after the same metavariable
assignment. If that assignment has cost one and reflexivity has cost zero,
the remaining cost is one, although there are two open obligations. The same
phenomenon occurs when a shared type or index choice makes several exact-local
steps possible. Merely requiring goals to be ``non-leaf'' before the assignment
does not prevent the counterexample.

Let $L_i(s)$ be an admissible lower bound obtained by relaxing all other
obligations while keeping goal $i$. Then
\[
 h_{\max}(s)=\max_i L_i(s)
\]
is admissible whenever the relaxed cost is no greater than the corresponding
joint cost. Summing the $L_i$ is justified only when costs have been partitioned
into non-overlapping resources, for example independent components satisfying
Proposition~\ref{prop:factor}. A more informative relaxation can retain a small
separator and compute a joint minimum over its finite abstract assignments.

There are two additional traps. First, rule-application cost is not necessarily
final-expression size: filling a shared hole can duplicate its term in the
fully expanded expression, while simplification can erase work. A bound for
one objective is not automatically a bound for the other. Second, zero-cost
rules can create infinitely many states of equal priority. Saturate a
terminating zero-cost normalization phase, or detect cycles in the zero-cost
state graph, before claiming that cost-ordered enumeration is fair.

\imp Begin with uniform-cost search under nonnegative integral costs and a
measured normalization phase. Add a heuristic only together with a small
counterexample suite and a statement of its relaxation. Maintain a separate
heuristic field for useful, unproved estimates. An estimate of likely subtree
size can guide exploration without being advertised as an A* lower bound.

\subsection{R1.3: can online cost learning coexist with defensible ranking?}\label{q:R1.3}
\ans Yes, if exploration priority and publication rank are separate. Online
learning may change what is discovered within a budget. That is a search
policy effect, not a logical soundness failure, and it should be visible in a
replayable record.

Freeze a publication objective $J_Q(p)$ when the query is created. It might be
a description-length cost plus explicit complexity penalties, as developed in
\qref{R8.1}. The scheduler may learn a different priority
\[
 P_t(s)=g(s)+\widehat h_t(s)-\lambda\widehat u_t(s),
\]
where $\widehat u_t$ estimates useful progress on observations. This expression
is a heuristic, not a lower bound unless independently justified. Learning
updates should occur at deterministic logical checkpoints, with stable tie
breaking and the same recorded observations.

\lit Probe establishes that partial solutions can supply useful just-in-time
guidance for bottom-up enumeration. It does not establish that adaptive
priorities preserve a particular user's preferred answer under a fixed time
budget~\cite{probe}. The architecture proposed here uses that guidance while
making the publication objective explicit.

Changing edge costs of an existing A* frontier invalidates its old optimality
argument. Either rebuild and consistently reprice the frontier under a fixed
new objective, or stop making an optimality claim and keep the learned score
solely as exploration guidance. A fair, deterministic baseline lane is useful:
it reserves progress for alternatives that a model currently undervalues.

Measure first certified result, best published rank at fixed work, and
held-out behavioral quality separately. Rewarding a candidate merely for
passing more training examples can favor overfitting. A change in result order
should be attributable to the frozen objective or to a recorded search-prefix
change, not to timing accidents in an asynchronous model response.

\subsection{R1.4: can wall-clock results be independent of machine speed?}\label{q:R1.4}
\ans Not unconditionally while also publishing every improvement reached
before the deadline. The appropriate guarantee is deterministic logical
prefixes, with explicit truncation at wall-clock cancellation.

\begin{proposition}[Deadline incompatibility]
Suppose a deterministic search trace discovers a better candidate after some
finite positive amount of additional work. If one machine reaches that point
before a fixed deadline and another does not, an anytime interface that
publishes its best discovered candidate returns different answers.
\end{proposition}
\begin{proof}
The faster execution has discovered the improvement and must publish it under
the stated interface; the slower execution has not. Returning the same answer
requires suppressing the improvement or obtaining work beyond the slower
execution's completed prefix. Neither is the stated interface.
\end{proof}

Startup calibration does not remove the issue: node costs vary with reduction,
allocation, instance synthesis, and proof search. Counting nodes can improve
comparability without accurately predicting time. A logical budget should
therefore include the expensive operations or use an explicitly defined
weighted ledger, while a wall-clock deadline remains a safety cancellation.

Publish a receipt containing the last completed logical checkpoint, exhausted
cost bands, uncompleted tasks, and the fixed candidate ordering. Parallel tasks
can be committed in deterministic batches rather than in completion order.
The UI may display earlier provisional results, but the stable checkpoint
result should be identified separately. Exact reproducibility across machines
then means replaying the same checkpoint and manifest, not pretending that the
machines perform the same computation in ten seconds.

The existing external ledger, which is not refunded by state rollback, is a
good foundation~\cite{implementation}. Extend it to persistent-frontier work
and proof tasks. Exclude timestamps, random scratch declaration names, and
memory addresses from canonical state and candidate identifiers.

\section{R2: carriers, accumulators, and invented state}\label{sec:r2}

\subsection{R2.1: can examples determine a minimal carrier type?}\label{q:R2.1}
\ans Only after specifying a bounded representation problem. Finite examples
can determine a smallest \emph{sample-consistent implementation in a finite
candidate language}; they do not identify the globally minimal carrier of an
unknown intended function.

Write a list implementation as
\begin{equation}\label{eq:carrier}
 h(x)=\phi\bigl(\foldr\,\delta\,z\,x\bigr),\qquad
 z:K,\quad\delta:A\to K\to K,\quad\phi:K\to T.
\end{equation}
Four objectives immediately differ: the cardinality of reachable states in
$K$; the size of the type expression denoting $K$; the description length of
$(K,z,\delta,\phi)$; and the operational cost of the resulting program. An
infinite carrier may have a tiny type expression. A tiny carrier type may move
all the difficult computation into its finisher.

\begin{example}[Why finite data does not identify the intended minimum]
For any finite set of list inputs, choose a word $w$ outside that set. The
constant-zero function and the function that returns one exactly on $w$ agree
on all those inputs when the observed outputs are zero. The former has a
one-state implementation; the latter cannot have a one-state implementation
with a deterministic finisher because it takes two distinct output values.
Both are total computable functions over a finite alphabet. Thus the data does
not determine even whether the intended function needs more than one state.
\end{example}

An unrestricted finisher makes carrier existence particularly uninformative:
choose $K=\List A$, $z=[]$, $\delta=\operatorname{cons}$, and $\phi=h$. The fold
reconstructs its input and the finisher performs the entire task. This is a
valid representation, but not an auxiliary-state discovery algorithm.

\paragraph{The bounded problem to implement.}
Specify finite grammars for carrier types, seeds, steps, and finishers; bounds
on syntax size, literals, providers, helper definitions, and universes; and a
fixed cost
\[
 J(K,z,\delta,\phi)=w_K|K|+w_z|z|+w_\delta|\delta|+w_\phi|\phi|.
\]
If the grammars describe a well-typed, total decidable fragment with an exact
validator, enumerate its candidates in nondecreasing cost, pruning only by
valid necessary conditions. A minimum satisfying the supplied observations
is then meaningful. For unrestricted Lean propositions, a candidate whose
proof search times out is \emph{unresolved}, so exhausting the enumerator does
not by itself establish nonexistence.

\begin{theorem}[Bounded synthesis guarantee]\label{thm:bounded}
Let $\mathcal P_b$ be a finite effectively enumerated set of well-typed candidate
implementations, let sample validation terminate exactly on every member, and
let every pruning predicate preserve all sample-consistent members. Exhaustive
cost-ordered search either returns a least-cost sample-consistent member or
correctly reports that none exists in $\mathcal P_b$.
\end{theorem}
\begin{proof}
Finiteness and termination of each validation imply termination of exhaustive
search. A valid pruning predicate removes no consistent member. Therefore the
first consistent member in cost order is minimal among consistent members;
if there is none, all members have been excluded by valid tests.
\end{proof}

The theorem applies to a finite reified DSL or a bounded collection with a
proved complete validator. It is not a theorem that arbitrary bounded Lean
metaprogram execution always finishes. A practical receipt should distinguish
\emph{enumeration exhausted} from \emph{all enumerated candidates resolved}.

\paragraph{Reverse is a cost example, not an existence obstruction.}
For ordinary lists,
\[
 \rev=\foldr\,(\lambda a\,r.\ r\concat[a])\,[] .
\]
This uses $K=\List A$. With linked-list append it repeatedly traverses the
accumulated suffix, giving quadratic work. The alternative
\[
 K=\List A\to\List A,\quad z=\operatorname{id},\quad
 \delta(a,k)=\lambda ys.\ k(a::ys),\quad\phi(k)=k([])
\]
uses one cons operation per element when its closures are applied. A restriction
on allowed step syntax may exclude the first program, but the fold
characterization alone cannot do so. Carrier synthesis should therefore
optimize the whole implementation and state its operational model.

\subsection{R2.2: what is the alphabet for a polymorphic function?}\label{q:R2.2}
\ans Abstract elements should be represented by symbolic tokens and their
permitted observations, not by a silently chosen finite concrete type.
Parametricity justifies this reduction only for a specified fragment.

For a total relationally parametric function
$h_A:\List A\to\List A$ that has no element-inspecting operations, renaming
elements along a function $f:A\to B$ gives the naturality equation
\[
 \operatorname{map}f(h_A(xs))=h_B(\operatorname{map}f(xs)).
\]
At a fixed input shape, a token model can record which input positions are
selected, duplicated, or reordered. It cannot manufacture a new abstract
$A$-value without a supplied source. Such a model is useful for proposing
carriers and rejecting impossible data flow. It still leaves infinitely many
list shapes and therefore does not establish an unrestricted finite-state
learning theorem.

If the context includes equality, order, hashing, or arbitrary callbacks,
the available observations change. An equality dictionary allows branches on
an equality pattern among tokens; an order dictionary allows order relations.
One then needs a symbolic state model with registers or constraints, rather
than merely ``one state per position.'' A \code{BEq} dictionary need not be
lawful, and its behavior is part of the actual program semantics. Replacing its
Boolean result with propositional equality requires an appropriate lawfulness
proof.

\lit Testing polymorphic properties has a substantial existing theory,
including reductions justified by parametricity~\cite{parametric-testing}.
That is evidence for exploiting a formally specified polymorphic fragment,
not permission to assume that every Lean term, every typeclass dictionary,
and every externally supplied provider satisfies the required relational
invariance.

\imp Give each provider a conservative effect signature: token-preserving,
shape-observing, equality-observing, order-observing, or opaque. Use positional
and equality-pattern abstractions to \emph{propose} states by default. Enable
hard parametricity-based pruning only for a certified provider fragment.
For indexed families, retain the index and the telescope in which a token is
typed; equal-looking tokens in different fibers are not interchangeable.

\subsection{R2.3: can fusion be run backwards, and what can automata prove?}\label{q:R2.3}
\ans There is a constructive derivation from a full specification with suitable
evidence, and there are complete bounded finite-state searches. Neither is a
general inference procedure from a handful of examples to a Lean carrier type.

\lit Gibbons, Hutton, and Altenkirch give the fold/unfold characterization;
Weber and Caldwell subsequently study constructive versions, including
formalization in Nuprl~\cite{fold-characterization,constructive-fold}.
Consequently, ``constructive fold characterization'' is not an entirely
unexplored question. Its premises concern the function and supporting
structure, rather than merely observed input--output pairs.

\begin{proposition}[A constructive sufficient condition for a list fold]
Let $h:\List A\to B$, and suppose a section $s:B\to\List A$ satisfies
$h(s(b))=b$. Assume cons congruence:
\[
 h(x)=h(y)\ \Longrightarrow\ h(a::x)=h(a::y).
\]
Then $h=\foldr\,\delta\,z$, where
$z=h([])$ and $\delta(a,b)=h(a::s(b))$.
\end{proposition}
\begin{proof}
The base equation holds by definition. For the inductive step,
$h(s(h(x)))=h(x)$, hence cons congruence gives
$h(a::s(h(x)))=h(a::x)$. Substituting the induction hypothesis for the folded
tail yields the desired equation.
\end{proof}

This is a genuine construction, but it consumes $h$, $s$, and a congruence
proof. It does not discover them from samples. For a polynomial functor $F$,
the analogous construction uses a section and an algebra formed by applying
$F(s)$ before the original constructor and $h$; a corresponding congruence
condition and induction principle justify it. This is a useful verified
transformation once a symbolic specification is available.

\paragraph{The precise residual-state model.}
For a full function $h:\List A\to T$, define
\begin{equation}\label{eq:residual}
 x\sim_h y\quad\Longleftrightarrow\quad
 \forall u:\List A,\ h(u\concat x)=h(u\concat y).
\end{equation}
The continuation is a \emph{prefix} because a right fold has already processed
the tail. This orientation avoids a common prefix/suffix ambiguity.

\begin{proposition}[Minimal semantic observer states]\label{prop:residual}
The quotient by $\sim_h$ supports a right-fold state machine with seed $[[]]$,
transition $\delta(a,[x])=[a::x]$, and output $\phi([x])=h(x)$. Every reachable
deterministic state representation satisfying \eqref{eq:carrier} maps onto
this quotient. In particular, a finite carrier needs at least as many reachable
states as the quotient has classes.
\end{proposition}
\begin{proof}
If $x\sim_h y$, then for every $u$,
$h(u\concat(a::x))=h((u\concat[a])\concat x)$ equals the corresponding expression
for $y$. Thus the transition is well-defined. The empty continuation makes
$\phi$ well-defined. Induction shows that folding $x$ produces its class.
For an arbitrary correct implementation, equal folded states for $x,y$ remain
equal after processing any prefix $u$, hence give equal final outputs; so
$x\sim_h y$. Mapping a reachable folded state to its residual class is
therefore well-defined and surjective.
\end{proof}

This is a semantic quotient, not necessarily finite, decidable, effectively
representable, or expressible by the chosen Lean type grammar. It distinguishes
what a state must remember without choosing a practical data structure.

\paragraph{What finite observations really imply.}
Create rows indexed by known tails $x$ and columns indexed by known prefixes
$u$. An entry is available only when $h(u\concat x)$ is observed or otherwise
proved. Put an incompatibility edge between two rows when some column has
\emph{two known different outputs}. Any clique of size $r$ certifies at least
$r$ distinct states. In contrast, two different partial rows may be compatible:
a row known only to be false in column zero and a row known only to be true in
column one can describe the same state. Counting distinct partial rows is
unsound as a lower-bound argument.

Even graph coloring is only a necessary constraint until transition
consistency is enforced: equal state labels for $x,y$ require equal labels for
$a::x,a::y$ whenever these are represented. A finite-state solver can introduce
state variables $q_x$, transitions $d_{a,q}$, and outputs $o_q$, with constraints
\[
 q_{a::x}=d_{a,q_x},\qquad o_{q_x}=h(x)
\]
for available observations, plus equality constraints for shared states.
Enumerating all transition/output tables at a fixed finite state bound gives
a decidable version of the problem. For the exact finite table language, an
exhaustive failure is a valid bounded negative result.

\lit Angluin's learning algorithm uses membership and equivalence queries to
learn regular languages~\cite{angluin}. A passive set of five to forty examples
is not such an oracle. Observation-table closure and consistency cannot be
assumed when the necessary entries are unknown.

A second correction concerns the proposal's count of constructible values.
The number of values \emph{found} at a search depth is normally a lower bound
on the available states, not an upper bound. Iterating a short step can also
produce values much larger than its own syntax. A state-count refutation is
valid for a carrier such as \code{Bool}, whose cardinality is proved, or for an
inductively closed abstract state set with a proved finite upper bound. It is
not justified by an incomplete value enumerator.

\imp Use incompatible observations to generate state requirements, propose
\code{Option}, products, and continuations, then synthesize the algebra and
finisher jointly. Where a full algebraic specification exists, try verified
fusion or tupling transformations. Otherwise retain fair bounded enumeration
as the completeness reference. The automaton is a constraint generator and,
in finite cases, a refuter; it is not a magical type decoder.

\subsection{R2.4: can motive enumeration cover primitive recursion efficiently?}\label{q:R2.4}
\ans A finite grammar can be enumerated completely up to its stated bound.
There is no reason to expect one small finite bound to cover every useful
primitive-recursive program with invented parameters. Efficiency requires
staging proposals without confusing their priority with completeness.

Use a dependent carrier grammar $K(\vec\imath)$ rather than only a closed type
$K$. Its first tiers can be the target family itself, functions from existing
parameter types, optional element state, and small products. Every proposal
must generate its seed, branch algebra, and finisher obligations at the same
time. A carrier that is cheap as a type but impossible to initialize is not a
promising partial answer.

Three examples explain useful priorities. An optional maximum state explicitly
represents ``no element seen'' and avoids inventing an abstract element for
an empty input. A Fibonacci state $(a,b)$ initialized to $(0,1)$ and updated to
$(b,a+b)$ turns the task into ordinary primitive recursion. A right fold
implementing a general left-associated reduction may need a continuation;
\code{Option A} alone should not be assumed sufficient for an arbitrary
non-associative combining operation.

For each total cost band, enumerate all well-scoped type expressions from the
finite grammar and all bounded implementations under them. Higher-priority
schema proposals may be tried first, but eventually process the entire band.
A beam or a fixed top-$k$ carrier list is an effective heuristic, not a complete
bounded search unless the omitted alternatives remain scheduled elsewhere.

Primitive recursion can be represented by iterated use of recursors, but this
observation alone does not solve the search problem: recursion depth, parameter
invention, auxiliary functions, and proof obligations need bounds too. State
completeness relative to the actual grammar, including its allowed recursor
nesting and parameter structure, rather than relative to all primitive
recursion in Lean.

\subsection{R2.5: do generalization and lemma-discovery heuristics transfer?}\label{q:R2.5}
\ans They transfer as proposal generators, especially when failure exposes a
specific missing invariant or auxiliary statistic. They do not supply a
complete or automatically trusted carrier learner.

A failed use of an induction hypothesis often indicates one of three changes:
an accumulator has a different value; the desired result contains extra
information absent from the recursive result; or a recursive result must be
combined through an unavailable lemma or helper. Respectively propose a
function-valued motive, a tupled state, or a helper with the residual hole's
typed observations. These are more targeted than enumerating every product
and function type at the outset.

\lit HipSpec combines theory exploration with inductive proving and uses
candidate equations generated and tested before proof attempts~\cite{hipspec}.
That gives a relevant architecture for conjecture generation and rejection.
It does not imply that tested equations are proofs, or that the same strategy
already invents arbitrary dependent carriers.

For Leant2, record a failure explanation: the induction hypothesis's type,
the required type, changed parameters, blocked equations, and the smallest
residual observations that distinguish failed branches. Generate a small
family of candidate invariants or helper types from that explanation.
Test them on known examples, then require native Lean proofs of the needed
branch conditions. Rank a proposed generalization by the obligations it
actually makes solvable, not just by a superficial resemblance between terms.

An especially useful pattern is strengthening
\[
 h(xs)\quad\rightsquigarrow\quad (h(xs),g(xs)),
\]
where $g$ supplies the missing statistic. The real obligation is closure:
both components on $a::xs$ must be computable from $a$ and the strengthened
state on $xs$. This closure test exposes the difference between inventing a
helpful statistic and merely appending another hard synthesis problem.

\imp Limit failure-driven invention by a deterministic budget and retain its
provenance. Reuse a helper only after its implementation and, when relevant,
its specification have been certified. Do not introduce a fresh helper axiom
and rely on the final program's type to reveal the unsound shortcut.

\section{R3: evaluation of partial programs}\label{sec:r3}

\subsection{R3.1: what is already known about higher-order and dependent examples?}\label{q:R3.1}
\ans Higher-order examples are not an entirely missing extension: Smyth's
implementation supports them. Dependent examples, however, introduce a
separate typing problem that cannot be solved by simply attaching ordinary
input--output pairs to a hole.

\lit Smyth's formal core and implementation have different scopes. The core
presentation uses first-order examples, while the implementation discussion
includes higher-order examples and polymorphism. Its live bidirectional
evaluation also addresses interdependent holes and examples that are not
trace-complete~\cite{smyth}. Therefore the proposal should not treat all those
capabilities as absent from the literature. The same source does not supply a
Lean-dependent implementation with Lean's conversion and transport semantics.

For Leant2, represent a hole closure by
\[
 (m,\Gamma_m,\rho,\tau[\rho],\mathcal O),
\]
where $m$ identifies the hole, $\Gamma_m$ is its declaration telescope, $\rho$
is the environment substitution, $\tau[\rho]$ is its instantiated expected
type, and $\mathcal O$ contains observations in that environment. The telescope
and substitution are not optional metadata. A hole of type
$\Pi n:\Nat,\,\operatorname{Vec}A\,n$ is observed in different fibers at
arguments zero and one; merging those result constraints as untyped values
would be meaningless.

For function-valued results, use application observations such as
$k(v_1)=w_1$ and $k(v_2)=w_2$. These are constraints on the function, not a claim
that a finite table completely describes it. A Church-encoded input can be
handled through the applications it performs on supplied eliminators, provided
closures and their environments remain explicit.

Backward propagation must also preserve alternatives. Unevaluating a required
output through a conditional may yield two disjunctive possibilities: a true
guard with a matching first branch, or a false guard with a matching second
branch. Choosing just one possibility is a heuristic restriction. The same
issue appears when an inverse operation is non-injective. Hard pruning must
consider every possibility in the declared bounded grammar or derive a
necessary condition valid for all completions.

\imp Stage the work. First extract independent observations while preserving
proofs relating them to the original contract. Next support typed first-order
hole closures. Then add higher-order application observations. Only after
those interfaces are stable add dependent fibers, casts, and generalized
motives. At every stage, an unsupported construct returns unknown rather than
being interpreted under an approximately matching type.

\subsection{R3.2: how should an assignment-driven normalizer survive backtracking?}\label{q:R3.2}
\ans Use a residual computation DAG whose cache entries carry validated read
sets and immutable assignment identities. Invalidation must distinguish two
sibling assignments even when a local version counter has been rolled back.

\lit Canonical-min describes suspended constraint continuations indexed by
the metavariables blocking them, with resumption after assignments. It is a
direct precedent for the requested control pattern, but not a drop-in
memoized normalizer for Lean expressions~\cite{canonical-min}. Adapton supplies
broader precedent for dependency-driven, demand-driven incremental
computation~\cite{adapton}. The implementation proposed here combines that
pattern with Lean state snapshots and typed residual observations.

A cache key should include the expression or closure identity, its environment
substitution, the relevant local telescope, the immutable Lean environment,
and the evaluator/transparency policy. Its read set records every consulted
metavariable assignment and universe assignment, including dependencies
reached through already assigned variables. Store the exact immutable
assignment-node identity, or a globally unique generation token, not a
counter that is restored with the branch.

For example, branch A may assign a hole to zero with version one and cache
that the hole equals one as false. After rollback, branch B may assign the
same hole to one, also with version one. Reusing the old result is wrong. A
branch-qualified immutable assignment identity prevents this ABA pattern.
Merely checking whether the hole is still ``assigned'' does not.

\begin{proposition}[Read-set validation]\label{prop:cache}
For a deterministic residual evaluator, suppose a cache entry records every
mutable value read during its computation and all other inputs are immutable.
If every recorded read has the same value in a new state, replaying the
computation produces the same result.
\end{proposition}
\begin{proof}
Assume replay first diverges at some operation. Earlier control flow and
intermediate results coincide. A pure operation on those results cannot cause
the first divergence. A mutable read cannot either, by the validated read set.
These exhaust the allowed operations, giving a contradiction.
\end{proof}

The theorem requires a complete read set. Tracking only holes that were
unassigned at the final stuck point misses earlier assigned values that
selected a branch. Also record dependencies on options, local instances,
and any environment extension that changes evaluation or inversion rules.

\imp Implement per-observation watchers before a full custom normalizer. An
assignment enqueues only affected observations, and each observation resumes
from a retained residual when its dependencies validate. Share closed
subexpressions globally within the frozen environment. Share open residuals
across branches only after read-set validation. Keep resource charges outside
rollback; a restored branch must not receive its spent evaluation budget back.

\subsection{R3.3: what is the cheapest sound fast-evaluation path?}\label{q:R3.3}
\ans Prefer an existing proof-producing evaluator and a narrow residual layer
before reimplementing the kernel. Also distinguish agreement with kernel
\emph{conversion} from agreement with the program's \emph{proved equations}.
The second is sufficient for certified refutation and can be substantially
more practical for well-founded recursion.

\lit Lean 4.34 documents \code{cbv} as call-by-value evaluation that uses
function and matcher equations and constructs propositional equality proofs.
\code{decide_cbv} applies it to decision procedures. Unlike native compiled
decision, this route does not trust the code generator; the documented proof
dependencies are standard Lean axioms, so the actual result must still be
audited against the selected profile. The same documentation explicitly
notes limitations on rewriting in dependent positions~\cite{lean-cbv}.
Lean4Lean's 2026 publication is important work on checking and verifying Lean,
but does not establish a complete, version-matched, verified replacement
normalizer for this project~\cite{lean4lean}.

This changes the engineering order. Use the native proof-producing machinery
for suitable closed observations and realized recursive definitions. Keep the
existing kernel reduction path where it works well. Add residual scheduling
and observation decomposition around these backends, rather than starting by
reproducing every kernel reduction feature.

A proof-producing evaluator may prove $e=v$ without $e$ and $v$ being
definitionally equal under the desired reduction strategy. That does not make
it an under-approximation of conversion. It makes it a semantics-preserving
rewriter with explicit evidence. If conversion itself is the intended
specification, choose a smaller fragment whose reductions are genuinely
kernel-convertible, and have the checker validate that property. Do not switch
between these claims in a performance argument.

For an open branch, a backend that reports false must support
\eqref{eq:refutation}. A useful implementation exposes three responses:
\begin{lstlisting}
proved(value, equalityProof, dependencies)
refuted(branchCertificate, dependencies)
unknown(residual, blockers, reason)
\end{lstlisting}
These are conceptual interfaces, not asserted Lean API declarations. A branch
certificate abstracts the holes into a well-typed telescope; it is not an
unclosed proof term containing unassigned metavariables.

Observation decomposition must be proved too. Splitting a conjunction is
straightforward. Translating a Boolean \code{BEq} test into propositional
identity is not justified without lawfulness. For a user-defined test,
retain its exact Boolean semantics unless a checked equivalence theorem
licenses the transformation.

\imp Initially let custom evaluation guide goal ordering and candidate
priority, while only existing checked paths prune. Then add certificate
production for a small fragment of constructors, applications, projections,
and specific recursor equations. Differential testing on generated programs
is a valuable bug detector, but even $10^5$ successful tests is not a proof of
conservativity. Include tests with open holes, delayed assignments, rollback,
universe constraints, and proof-dependent types; closed-term agreement alone
misses the hardest failure modes.

\subsection{R3.4: does top-down angelic recursion admit correctness guarantees?}\label{q:R3.4}
\ans Yes under explicit boundedness and completeness assumptions. Top-down
organization is not itself an obstacle. The difficult parts are preserving
all possible recursive results, enforcing their eventual consistency, and
not confusing an incomplete oracle solver's failure with impossibility.

A recursive call on an unknown smaller input can be interpreted by a symbolic
result variable. Calls with the same complete argument tuple should share the
same variable when their deterministic semantics requires it. The tuple must
include accumulators and relevant environment values. Omitting that correlation
may be a safe but very weak over-approximation; imposing a guessed correlation
between unequal calls can incorrectly prune.

Known observations can constrain the oracle when the contract actually entails
them. If $C(f)$ implies $f(v)=w$, the oracle may use $w$ for that call while
searching for a function satisfying $C$. A contradiction obtained this way
refutes \emph{contract-satisfying} completions. It need not show that every
arbitrary completion evaluates the observation to false. The certificate should
therefore retain the contract implication explicitly.

\lit Burst supplies angelic synthesis with strengthening and backtracking.
Its completeness statements depend on the behavior of the angelic synthesis
subroutine and the validity of failure information; they are not an
unconditional termination theorem for arbitrary recursive Lean synthesis~\cite{burst}.

\begin{proposition}[A sufficient bounded top-down setting]
Fix a finite candidate-body grammar, a finite admissible oracle-call/value
space, a well-founded recursive-call relation, an exact terminating solver for
the resulting constraints, and a fair traversal that resolves each finite case
once and stops on exhaustion. If refinement excludes only
infeasible interpretations or invalid completed candidates, top-down angelic
search terminates and is complete for that bounded candidate problem.
\end{proposition}
\begin{proof}
There are finitely many candidate bodies and finitely many oracle assignments
for each. Exact constraint solving and well-founded evaluation resolve each
case. Valid refinement removes no successful body with its genuine recursive
results. Fair traversal therefore either reaches such a body or exhausts all
possibilities.
\end{proof}

Finite syntax alone does not guarantee the other assumptions. Recursive calls
may involve large values or new arguments, and solver timeouts remain unknown.
For an initial implementation, use angelic assumptions only to prioritize
search, then actually evaluate and certify each completed body. Promote angelic
infeasibility to hard pruning only when the relevant finite relaxation has a
checked or otherwise established complete refuter.

\section{R4: dependent motives, generalization, and transport}\label{sec:r4}

\subsection{R4.1: can a useful motive class be enumerated completely?}\label{q:R4.1}
\ans Yes for a precisely specified finite syntactic class. No, abstraction of
maximal index occurrences alone is not closed under the inventions required
for accumulators and strengthened results.

Suppose a goal contains finitely many selected maximal occurrences of index
terms. Every subset determines a candidate abstraction, after scope and typing
checks. Enumerating all subsets terminates and is complete for \emph{that
construction}. Most candidates may fail to typecheck; that is an efficiency
issue, not a reason to infer completeness for a broader class.

An accumulator parameter $S$ may not occur in the goal's indices at all.
Similarly, a strengthened result $T\times U$ may introduce $U$ for the first
time. Neither can be produced solely by replacing existing index occurrences.
Extend the class with bounded dependent carrier templates and
\emph{dependency-closed generalization sets}. If a local is reverted, all later
locals whose types or values depend on it must be handled in telescope order.

For vector map, a natural motive is
\[
 M(n,x)=\operatorname{Vec}B\,n.
\]
The function argument $A\to B$ may remain a fixed parameter. In an accumulating
variant, a useful motive might instead be
\[
 M(n,x)=\Pi s:S,\,T(n,x,s).
\]
The second is not merely another subset of the first goal's index occurrences.
A unified motive grammar makes the relationship to carrier invention explicit.

Compound and repeated indices need special treatment. A major premise such as
$\operatorname{Vec}A(n+m)$ cannot always be handled by pretending that $n+m$ is
an independent local variable and forgetting its defining equation. Generalize
the compound expression to a fresh index, retain the equality and its dependent
consequences, and let native elimination construct the resulting obligations.

\imp Enumerate motives in stages: native inferred motive; native motive after
minimal dependency-closed reversion; bounded extra parameter or product
strengthening; finally a wider syntactic abstraction grammar. Record the
stage and bound. Use a fair fallback to preserve relative completeness;
otherwise explicitly report a heuristic search rather than ``no motive
exists.''

\subsection{R4.2: how can transports be handled without destroying identity?}\label{q:R4.2}
\ans Treat transports as typed edges with checked equality evidence. Normalize
only through justified rules, and distinguish proof-irrelevant evidence from
the type-changing operation that evidence enables.

The proposal's example \code{Vec A (n + 0)} versus \code{Vec A n} does not
require a propositional cast: Lean reduces addition on its second argument.
\code{Vec A (0 + n)} is a suitable replacement when $n$ is a variable; proving
$0+n=n$ is then relevant~\cite{lean-recursion}.

For a family $F:I\to\operatorname{Sort}u$ and $e:i=j$, transport maps $F(i)$
to $F(j)$. Proof irrelevance can identify different proofs of the same
proposition; it does not identify the two endpoint types, erase the family,
or justify dropping the transport from a term that must typecheck. A printed
presentation may hide a cast, but the replayable artifact must preserve it.

A practical representation keeps an ordinary Lean expression plus an auxiliary
transport record containing source type, target type, equality proof, and
normalization provenance. For arithmetic indices, orient checked equations
into a chosen normal form and reuse a cached equality certificate. Eliminate
identity casts and compose compatible paths when a proved rule permits it.
Do not enumerate arbitrary equality proofs: synthesize the equation once,
then let proof irrelevance and typed proof caching remove redundant evidence.

There is no general cheap canonical quotient of all dependent Lean programs
by propositional equality. Setoids change the specification and require their
own compatibility proofs; heterogeneous equality changes the form of the
obligations but does not automatically construct the desired homogeneous
inhabitant. Both can be useful in specialized intermediate theories, not as a
universal shortcut around native dependent typing.

\imp Bound nontrivial transport steps and assign them a nonzero structural
penalty unless a certified normalization removes them. The proposal to make
all casts zero-cost creates avoidable cycles and makes a size objective
misleading. For display, provide a readable expression and an exact checked
expression; do not use the readable form as the identity key.

\subsection{R4.3: what unfolding policy should search use?}\label{q:R4.3}
\ans Use staged, demand-driven unfolding with an explicit policy in the state
and cache keys. Neither global normalization nor a single maximal transparency
setting is a convincing default for large dependent search.

\lit Canonical's 2025 account describes unfolding certain aliases that expose
function types, reducible definitions, and typeclass-related structure as part
of translation into its own representation~\cite{canonical}. That is concrete
evidence for selective unfolding, not proof that the same translation policy
is optimal inside Lean's metaprogramming engine. The \code{sauto} publication
establishes a relevant dependent type-inhabitation search approach~\cite{sauto};
the precise current implementation-level unfolding policy was not verified in
this investigation, so it should not be represented as a confirmed benchmark
result or copied by attribution alone.

The proposed native policy has three stages. First perform cheap weak-head
normalization sufficient to expose a function, constructor, or indexed-family
head. Second, when matching is blocked on a reducible named definition, unfold
only the demanded head under a fixed budget. Third, permit a bounded stronger
normalization or theorem-based rewrite for a specific failed match. An opaque
constant should remain opaque unless the permitted interface supplies a
checked equation or another explicit reason to expose it.

Record why unfolding was requested: exposing a binder, determining a recursor,
matching an index, or evaluating an observation. These requests have different
cost/benefit profiles. Cache a result only under the same environment,
transparency options, local context, and assignment dependencies. Do not let a
failed low-transparency match become a permanent negative fact about the term.

\imp Compare three policies on the same dependent cases: fixed weak-head
normalization, staged demand-driven unfolding, and bounded up-front
normalization. Measure successful matches, reduction work, and missed
candidates at fixed logical budgets. Treat broader comparative claims about
Canonical or \code{sauto} as experiments to perform, not as existing evidence
of superiority.

\subsection{R4.4: is native motive computation available as a library?}\label{q:R4.4}
\ans Yes. The inspected Lean 4.34 source exposes useful components, but they
have concrete preconditions and are not a universal motive-enumeration API.

The public native operation includes
\begin{lstlisting}
Lean.MVarId.induction
  (mvarId : MVarId) (majorFVarId : FVarId)
  (recursorName : Name) (givenNames : Array AltVarNames := #[]) :
  MetaM (Array InductionSubgoal)
\end{lstlisting}
The result includes each subgoal, its fields, and a free-variable substitution.
The same file exposes \code{getMajorTypeIndices} and
\code{mkRecursorAppPrefix}. The latter builds the recursor prefix and computes
a motive by abstracting the major premise and indices. The index-extraction
path checks that indices are suitable distinct variables and reports specific
failures when they are not. Native induction reverts dependencies and
reintroduces the major premise and indices~\cite{lean-induction-source}.

Wrap these operations in a version-pinned adapter. Each attempt runs in a
fresh saved state, creates no globally visible declaration, and returns a
structured outcome: success with subgoal substitutions; compound-index
precondition failure; target-sort/elimination restriction; failed unification;
or resource exhaustion. Preserve interrupts and do not classify a timeout as
an impossible motive.

For a compound index, first use an appropriate native generalization path
rather than invoking the low-level operation with violated preconditions.
After native induction, retain its substitutions instead of assuming that
original free-variable identifiers still denote the branch locals. This is
especially important when attaching residual examples to new branch holes.

\imp The first dependent milestone should be vector map and a small family of
indexed constructors using this adapter, not a replacement recursor compiler.
Benchmark the cost of both successful and failed attempts. ``Cheap failure''
is a performance target to measure; it is not guaranteed merely because the
function is callable from \code{MetaM}.

\section{R5: recursive bodies and checked realization}\label{sec:r5}

\subsection{R5.1: does a capability interpreter agree with the realized definition?}\label{q:R5.1}
\ans Agreement must be proved for the chosen body language and realization.
For well-founded recursion, expect a propositional equation rather than
assuming a convenient definitional reduction. Native equation-based evaluation
is the first implementation to try.

A dependent capability has the form
\[
 \mathit{rec}_x:(y:\alpha)\to R(y,x)\to C(y),
\]
where $R$ is well-founded and $x$ contains every argument relevant to recursive
state. A body functional has type
\[
 F:(x:\alpha)\to\bigl((y:\alpha)\to R(y,x)\to C(y)\bigr)\to C(x).
\]
An accumulator is not an invisible interpreter variable: either it belongs in
$x$, or the result family explicitly quantifies over it. This prevents oracle
memoization and termination checking from silently treating unequal calls as
the same call.

Compile a completed, typed body using a chosen native realization, obtaining
$f$ and its checked equation
\begin{equation}\label{eq:rec-equation}
 f(x)=F\bigl(x,\lambda y\,h.\ f(y)\bigr).
\end{equation}
The exact equation may include the elaborator's administrative arguments and
casts. The native recursion documentation explicitly distinguishes convenient
structural reduction from the potentially non-definitional equations of
well-founded recursion~\cite{lean-recursion}.

A direct interpreter can be related to $f$ by well-founded induction. At each
recursive call, apply the induction hypothesis to the smaller argument and
use congruence for the body constructs. This proves agreement if the body
interpreter preserves the semantics of every construct and the realization
has equation~\eqref{eq:rec-equation}. Neither condition follows merely from
using similar-looking source code.

A fuelled interpreter needs a stronger statement:
\[
 \operatorname{run}(n,x)=\operatorname{some}(v)\ \Longrightarrow\ f(x)=v.
\]
Its failure to produce a value with fuel $n$ does not imply that $f(x)$ has a
wrong value. The fuel bound may be inadequate even for a terminating function.
Return unknown, retain the residual, and charge the work.

\imp Reuse the proof-producing evaluation path of \qref{R3.3} on realized
candidates and their equations. For open-body exploration, interpret only a
small typed capability language, initially without hard pruning from
uncertified results. Realize promising completed bodies once, cache their
checked equations, and certify the contract. Never insert a global temporary
axiom naming the recursive function and later forget to remove it from the
accepted dependency closure.

There is also a distinction between mathematical acceptance and deployment:
a candidate can kernel-check yet encounter compilation or binding problems in
the front end. Since the proposal records such an example for a synthesized
recursor term, publication should include a separate executable-binding check
when the requested output is an executable definition~\cite{proposal-evidence}.

\subsection{R5.2: what changes when angelic execution is top-down?}\label{q:R5.2}
\ans The semantics of the relaxation can remain the same; the representation
of alternatives and the way constraints propagate change. Top-down search
trades the compact bottom-up representation of many terms for goal-directed
refinement of a typed sketch.

A top-down node contains a body skeleton, its holes, recursive-call variables,
and constraints derived from observations. Filling a constructor hole can
immediately propagate a required result to its fields. Filling a guard can
select which branch constraints apply. Filling a recursive call's arguments
can identify a previously symbolic call with an observed input or another
call. These changes naturally fit the dependency watchers of R1 and R3.

The challenge is retaining correlated alternatives. If one observation demands
that a recursive result be either zero or one, and another demands that the
same call be positive, prematurely choosing zero loses the feasible option.
Store the alternatives symbolically, or branch over all of them within the
finite bound. A solver that only tries one satisfying angelic assignment can
still be a useful proposer, but its later failure does not refute the body.

Burst's bottom-up method uses finite-tree-automaton representations and a
particular refinement organization~\cite{burst}. A top-down implementation
cannot inherit the space savings or the proof assumptions simply by using the
word ``angelic.'' It needs its own constraint representation and the bounded
correctness conditions stated in \qref{R3.4}.

\imp Start with recursive calls at already observed inputs, then symbolic
unknown results used for ranking, and finally exact bounded oracle constraints
that may support pruning. Keep actual closed execution and kernel-certified
contracts as the acceptance path. This gives useful intermediate milestones
without making the entire search rewrite depend on a general angelic solver.

\subsection{R5.3: is there a small canonical recursion basis with measured coverage?}\label{q:R5.3}
\ans There is a small useful engineering basis, but no coverage percentage for
Leant2 follows from that fact. A universal semantic basis and an efficient
search basis are different objectives.

A proposed staged basis is: structural recursion on one major premise;
structural recursion with generalized parameters or tupled results; matching
additional arguments inside that recursion; bounded course-of-values access;
and well-founded recursion under a small measure grammar. Mutual or nested
families can initially use native elaboration when an appropriate body sketch
has been found, rather than expanding every generated recursor in the inner
loop.

Several textbook examples need less machinery than their surface form
suggests. Zip can recurse on one list while matching the other. Fibonacci can
use the pair state in R2 rather than directly invoking two sub-subterm calls.
Tree flattening can fold into difference lists. Merge can use an explicit
measure such as the sum of the two list lengths, whereas a search that insists
on a fixed single structural argument may have difficulty. These examples
motivate the tiers; they do not establish empirical coverage on a corpus.

A body capability should encode exactly which calls are available. A
``course-of-values depth two'' restriction is a syntactic capability bound,
not a generic theorem about all course-of-values definitions. A measure grammar
must also bound arithmetic expressions, lexicographic tuples, and proof search
for decrease. Otherwise a finite list of scheme names hides an unbounded
secondary synthesis problem.

\imp Classify reference programs before measuring synthesis: minimum sufficient
scheme tier, whether a carrier is supplied, whether helpers are supplied,
whether examples are trace-complete, and whether a universal proof is required.
Run both provider-enabled and provider-withheld tracks. Report representation
coverage separately from search success; a program may be expressible in tier
one yet not found within the budget. No numerical coverage claim is made here.

\section{R6: contracts, induction, and proof debt}\label{sec:r6}

\subsection{R6.1: how should value holes and proof holes be interleaved?}\label{q:R6.1}
\ans Use dependency-triggered proof tasks with a cheap eager tier and a bounded
expensive tier. Waiting for the entire program is too late, but declaring every
proof goal meaningless while it contains holes is too strong.

A proof obligation can simplify or produce useful equalities before all its
value dependencies are known. A proposition such as $x=x$ needs no value of
$x$; constructor discrimination can refute some partial branches immediately.
Other obligations genuinely block on a value choice. Represent that distinction
with dependencies rather than a blanket ``proof last'' rule.

For each proof task, retain its exact context, target, relevant value holes,
attempted tactic tier, and certified consequences. Cheap normalization,
reflexivity, and explicit contradiction checks run whenever a dependency
changes. More expensive tactics run only after the relevant branch values are
stable enough or the candidate has survived inexpensive tests. A timeout
suspends the task; it does not certify that the chosen value is unprovable.

\lit Synquid demonstrates synthesis directed by polymorphic refinements with
automated constraint solving~\cite{synquid}. Its logical fragment and solver
assumptions are not those of arbitrary Lean propositions. The analogue proposed
here is an explicit finite abstraction of selected predicates, with all
consequences justified by Lean proofs rather than silently assuming a complete
SMT interpretation of \code{Prop}.

A useful optional layer maintains abstract predicates such as length bounds,
non-emptiness, or membership facts. Each transfer rule has a checked lemma.
The abstract domain can be incomplete: failure to infer a property merely
leaves an obligation. This supplies a Lean-native refinement discipline
without requiring a complete decision procedure for every contract.

For subtype targets, synthesize a branch value and its local proof together.
A recursive result then carries the induction invariant needed by the next
branch. Reuse that proof rather than reconstructing a global theorem about an
opaque recursor expression after synthesis has finished.

\imp Make proof debt a first-class queue over the same dependency graph as
value search. Wake a task only when its blockers change or its allocated tier
increases. Record whether a proof was attempted before or after the relevant
value became stable; this exposes wasted repeated work and delayed failures
in a way that a single aggregate proof timer cannot.

\subsection{R6.2: which properties have a recognizable induction-along-the-program proof?}\label{q:R6.2}
\ans A useful sufficient class can be recognized through generated verification
conditions and a certified leaf solver. There is no need to decide membership
in the much broader class of all properties for which a clever induction,
generalization, or auxiliary lemma exists.

For a fold $s(xs)=\foldr\,\delta\,z\,xs$, let $I(xs,s)$ be a supplied or
synthesized invariant. Generate the obligations
\begin{align}
 &I([],z),\label{eq:vcbase}\\
 &\forall a,xs,r,\ I(xs,r)\to I(a::xs,\delta(a,r)),\label{eq:vcstep}\\
 &\forall xs,r,\ I(xs,r)\to C(xs,\phi(r)).\label{eq:vcfinish}
\end{align}
The input $xs$ remains in the invariant, so this scheme expresses relationships
between the input and the accumulated state, not only properties of the state.

\begin{theorem}[Local invariant certification]\label{thm:invariant}
Proofs of \eqref{eq:vcbase}--\eqref{eq:vcfinish} imply
$\forall xs,\ C(xs,\phi(\foldr\,\delta\,z\,xs))$.
\end{theorem}
\begin{proof}
Induction on $xs$ proves $I(xs,\foldr\,\delta\,z\,xs)$. The empty case is
\eqref{eq:vcbase}; the cons case follows from the induction hypothesis and
\eqref{eq:vcstep}. Apply \eqref{eq:vcfinish} to obtain the contract.
\end{proof}

This theorem provides a concrete proof template. The same construction can
be internalized by a motive
\[
 xs\longmapsto\{r:K\mid I(xs,r)\},
\]
so that each recursive branch returns a value and the required evidence.
Projecting the value yields the program while the proof projection establishes
its invariant.

For map with element function $f$, choose
$I(xs,ys)\equiv |ys|=|xs|$. The cons verification condition reduces to equality
of successor lengths after using the induction hypothesis. For a numerical
accumulator, an affine invariant may reduce the step to linear arithmetic.
For reverse or sorting, an adequate invariant may require append, permutation,
or ordering lemmas; those are additional inputs to the proof template, not
consequences of the word ``structural.''

Define a recognizer for a sufficient fragment by allowing finitely many
invariant templates, a terminating oriented rewrite system with proved rules,
and a checked decision procedure for the residual arithmetic or constructor
formulae. Acceptance by this recognizer produces a proof. Failure means only
that the selected template/solver combination did not certify the property.
Even when every leaf lies in a decidable theory, discovering the right
invariant can remain the difficult synthesis task.

\imp Expose a separate outcome \emph{template applicable but invariant missing}.
This can trigger R2's generalization or helper proposals. Do not repeatedly
spend the largest proof budget on a monolithic goal whose branch obligations
show that the current invariant cannot connect the recursive states.

\subsection{R6.3: what Lean induction automation is usable now?}\label{q:R6.3}
\ans Use native induction and functional-induction infrastructure to choose a
small set of explicit proof skeletons, then apply a bounded leaf portfolio.
Do not assume that a general-purpose tactic automatically discovers all useful
induction variables, generalizations, and lemmas.

The inspected native induction interface already supplies recursor-based
subgoals and substitutions, as explained in \qref{R4.4}. Lean's reference
manual also documents functional induction associated with recursive
definitions~\cite{lean-recursion}. These are appropriate infrastructure for a
proof that follows the synthesized program's recursion. The choice of an
adequate invariant remains the synthesizer's responsibility.

A practical core-first sequence is: introduce quantified inputs; select the
major premise or the realized function's induction principle; apply the known
invariant/generalization template; simplify with the program's checked
equations and local induction hypotheses; then try the relevant arithmetic or
congruence solver. Guard every call by explicit work limits and preserve the
state on failure. \code{grind} belongs in the leaf portfolio, not in the design
as an unspecified promise of automatic lemma discovery.

\lit Canonical is directly relevant to automated dependent inhabitation and
recursor-based reasoning~\cite{canonical}. LeanHammer's current work supplies
learned premise selection and proof reconstruction infrastructure, including
a 2026 conference publication~\cite{leanhammer}. These are distinct resources:
neither a theorem-retrieval accuracy figure nor success on ordinary proof
corpora establishes automatic induction coverage for newly synthesized
programs.

External hammers and cloud retrieval should remain optional extensions, not a
hidden dependency of the core-only path. Their availability, latency, and
axiom effects must be included in the receipt. A synthesized proof term that
reconstructs inside the permitted profile is valuable; a timeout or failed
external reconstruction is simply unresolved.

\imp Build a regression suite with hand-classified proofs: one induction plus
simplification; induction plus arithmetic; generalization required; helper
lemma required; stronger invariant required. This is more informative than
one overall solved count and permits targeted expansion of the automation
without changing the semantics of an accepted result.

\subsection{R6.4: is there a principled policy for proof debt?}\label{q:R6.4}
\ans There is a principled decision-theoretic objective, but no reason to claim
that ordinary independent-arm bandit theory yields an optimal policy for this
coupled search. Proof tasks share candidates, lemmas, contexts, and setup work.

A reasonable heuristic score is
\[
 \operatorname{priority}(j)=
 \frac{\widehat{\Pr}(\text{useful certification or refutation}\mid j)
       \cdot\widehat{\Delta U}_j}
      {\widehat{\operatorname{work}}_j+\widehat{\operatorname{switch}}_j+\varepsilon}.
\]
The utility can reward a better certified candidate, the rejection of a large
branch, or a reusable local lemma. Estimates should be learned or calibrated
from held-out traces. This formula is a proposed heuristic, not a theorem of
optimal scheduling.

Use deterministic epochs and minimum service quotas. Each epoch assigns work
to value expansion, residual evaluation, and proof tasks. Within the proof
queue, prefer stable local obligations over large global ones, and reuse a
warm context when doing so saves real setup work. Escalate a task from cheap
to expensive tactics only when its previous tier has completed and a relevant
reason exists to spend more. A fair quota prevents uncertain but essential
proof obligations from starving forever.

The key scheduling unit is a \emph{task continuation}, not necessarily an
entire tactic invocation that restarts from scratch. Some native tactics may
not expose resumable internals; account for that honestly and cache their
completed preprocessing where possible. Do not implement ``500 milliseconds
per proof'' and then claim deterministic results. Prefer logical limits tied
to the pinned evaluator/tactic version, with wall-clock cancellation as a
separate status.

\imp Log the number of obligations created, woken, discharged, refuted,
retried without changed dependencies, and timed out. Compare a fixed staged
policy with learned prioritization under equal total work. Include total
candidate quality, not just tactic success rates: cheaply proving properties
of low-value candidates can otherwise look artificially good.

\section{R7: library retrieval with dependency and type classes}\label{sec:r7}

\subsection{R7.1: what abstraction is sound and useful for dependent signatures?}\label{q:R7.1}
\ans Use a progressively refined typed hypergraph, retaining constructor and
index shape and treating unknown information as an over-approximation. A graph
of isolated type heads is a useful first ranking feature, but not a complete
synthesis semantics.

A component $c:A_1\to\cdots\to A_k\to B$ needs all of its arguments, not merely
one edge from an available input head to $B$. Use a hyperedge whose premises
record the argument abstractions and relevant shared type variables. Ordinary
Lean contexts permit reusing variables, so the abstraction must also model
contraction rather than accidentally treating every value as a linear token.
Lambda introduction, locally available functions, constructors, and permitted
elimination rules need their own abstract transitions if the abstraction is
to justify negative reachability.

An initial abstract type can retain: the outer constructor; the arrangement
of type arguments to bounded depth; an equality partition of repeated type
variables; coarse index forms such as zero/successor/unknown; and necessary
universe relationships. Erase an unresolved index to an explicit top element,
not to a guessed constant. Preserve distinctions that explain common failed
matches: \code{Vec A n} and \code{Vec A (n+1)}, or a function whose input and
output share the same type parameter versus one with unrelated parameters.

When a proposed composition fails concrete Lean checking, refine only the
abstraction responsible for the spurious connection. A term-index mismatch can
split an index class; a missing dictionary can add a class precondition; a
universe contradiction can distinguish level patterns. A timeout supplies no
such counterexample and must not be used as a refinement fact.

\begin{proposition}[Safe abstract non-reachability]
Suppose every concrete construction rule in a selected grammar is simulated
by an abstract transition, and every concrete initial resource maps to an
abstract initial resource. If the abstract target is unreachable, no concrete
term in that grammar reaches the target.
\end{proposition}
\begin{proof}
Induct on a concrete term's construction derivation. Initial resources are
represented by assumption. Every subsequent construction maps to an abstract
transition whose premises are reachable by the induction hypothesis. Thus any
concrete solution would induce an abstract path to the target, contradicting
non-reachability.
\end{proof}

The simulation assumption is the whole issue. A depth-three library path
filter that omits recursors, lambdas, casts, or newly invented helpers may still
rank providers effectively, but it cannot establish nonexistence for a search
that uses those constructs. Likewise, erasing universes is safe only when it
enlarges the possible matches rather than introducing a restrictive default.

\imp Initially use abstract reachability to prioritize a local provider set,
with a wider deterministic fallback. Promote negative results to hard pruning
only for a documented finite component sublanguage with a simulation argument.
Measure selectivity and concrete mismatch causes; a universal type-variable
node that connects everything signals a need for refinement, not a reason to
silently discard inconvenient polymorphic functions.

\subsection{R7.2: has type-guided refinement already handled type classes?}\label{q:R7.2}
\ans Yes. Hoogle+ explicitly represents typeclass constraints by dictionaries.
The remaining research problem is extending the approach to Lean's dependent
and universe-sensitive setting, not inventing dictionary treatment from
scratch.

\lit Section 5.1 of Hoogle+'s paper describes dictionary types, extra dictionary
arguments for constrained methods, and instance functions that construct
result dictionaries from prerequisite dictionaries. It also states limits on
higher-kinded type variables~\cite{hoogle}. This directly answers the proposal's
question about whether type-guided abstraction refinement has been tried with
type classes, while leaving important differences from Lean unresolved.

For example, abstract a constrained operation schematically as
\[
 \mathit{eqMethod}:\operatorname{EqD}(A)\to A\to A\to\Bool.
\]
An instance for lists is another component
\[
 \mathit{eqList}:\operatorname{EqD}(A)\to
                    \operatorname{EqD}(\List A).
\]
The reachability problem now includes obtaining dictionaries. In Lean the
actual dictionary terms, local instances, parameters, and output-parameter
behavior must be interpreted by native instance synthesis rather than by an
independent approximation that silently becomes authoritative.

The proposed cache key ``class and type'' is insufficient in general. Add the
local instance context, environment fingerprint, universe instantiation,
relevant options, and dependencies on unresolved metavariables. An instance
search that is stuck on an output parameter is not a cached global failure.
A successful cached dictionary should be reinstantiated and typechecked in the
new context; proof-bearing or data-bearing class fields remain part of the
actual term.

\imp Index constrained components with explicit dictionary prerequisites,
then resolve concrete obligations lazily once their dependent parameters are
sufficiently known. Cache positive dictionaries aggressively under validated
keys. Keep negative instance information local to the complete, fixed search
state and bound that justified it. Benchmark instance-search work separately
from ordinary provider matching, since an apparently small provider set can
still open a large instance search.

\subsection{R7.3: can premise selection consume data goals and examples?}\label{q:R7.3}
\ans The current Lean interface can be called on a data goal and has a hint
field, but this does not imply that an installed selector understands examples
or has been trained for program-component relevance. A dedicated adapter and
training/evaluation corpus are needed.

The inspected Lean 4.34 API has
\begin{lstlisting}
Selector := MVarId -> Config -> MetaM (Array Suggestion)
\end{lstlisting}
with configuration fields for a suggestion limit, caller, filter, and an
optional arbitrary hint. Its type does not restrict the goal to \code{Prop};
the hint is explicitly not guaranteed to be used by a selector~\cite{lean-selector-source}.
Consequently an example summary can be passed through an adapter, but claiming
that this automatically supplies behavior-aware retrieval would be unwarranted.

\lit The current cloud premise-selector documentation describes goal-state
and premise embeddings, and warns about substantial cold-start costs. It is
therefore not a dependable unbudgeted operation inside a short core-only
query~\cite{selector-readme}. Its documented interface and the LeanHammer
research motivate a reusable retrieval backend, not an established score for
carrier or data-program synthesis.

Construct a program-retrieval corpus from executable definitions in Std,
Mathlib, and appropriately licensed libraries. Hide the target body; collect
its signature, permitted environment, generated or specified observations,
and the dependency components actually used by its implementation. Multiple
implementations may be valid, so dependency prediction is a training signal,
not a unique ground truth for intended behavior.

Leakage control is essential. Remove the target declaration itself, aliases,
and providers that transitively expose its body when evaluating genuine
synthesis. Split related function families and dependency strongly connected
components between training and test where possible. Otherwise a held-out
``definition'' may be recoverable by an imported synonym. Separate
provider-enabled API search from provider-withheld algorithm invention.

For behavior features, compute summaries such as length change, shape change,
order preservation on observed inputs, and use of supplied callbacks. These
are relevance signals, not logical filters unless certified. A component that
fails to equal the final answer may still be an essential intermediate step.

\imp Keep the existing selector interface for proposition goals and implement
a data-goal wrapper that combines type compatibility, abstract reachability,
and behavioral features. Cache indexes per environment. The default lane
should remain local and deterministic; optional remote retrieval must use a
recorded result manifest and a separately charged latency budget.

\section{R8: ranking, equivalence, and ambiguity}\label{sec:r8}

\subsection{R8.1: what formal objective should rank inhabitants?}\label{q:R8.1}
\ans A frozen minimum-description-length objective is defensible as an explicit
prior, not as a universal definition of user intent. Correctness is a hard
constraint; program quality is a separate, testable preference model.

Choose a prefix-decodable or otherwise explicitly weighted grammar and define
\[
 J_Q(p)=L_{\mathcal G_Q}(p)+\lambda_1\operatorname{casts}(p)
       +\lambda_2\operatorname{helpers}(p)+\lambda_3\operatorname{risk}(p).
\]
Here $L$ counts the complete representation, including carrier, finisher, and
helper definitions, not just the visible outer term. A library call can be
cheap if the grammar assigns it a short code, but this is a stated prior over
the available library, not an objective fact that the called implementation
is computationally simple. Any runtime or space cost model should be an
additional explicit objective with its own evidence.

For finite examples $E$, an MDL-style interpretation is
\[
 p^*\in\arg\min_{p:\,p\models E} L_{\mathcal G_Q}(p).
\]
For a universal contract, replace the sample constraint with the existence of
a checked proof. Neither formulation says that the user intended the smallest
consistent program. It does make the ranking criterion reproducible and
inspectable.

\lit Learned ranking has empirical precedent in practical programming by
example, including the real-time ranking work of Natarajan and
colleagues~\cite{ranking-study}. User interaction for disambiguating consistent
programs has also been studied~\cite{pbe-interaction}. These are reasons to
evaluate ranking against user intent and ambiguity; they do not establish that
MDL or Leant2's unused-input heuristic is optimal for dependent proof-assistant
programs.

Unused-input penalties are especially delicate. A correct constant function
or projection may intentionally ignore arguments. Syntactic occurrence is not
information flow: $x-x$ mentions $x$ but, under ordinary arithmetic semantics,
is constant. Conversely, a dependency hidden in a callback or dictionary may
be operationally important. Use such penalties as weak, documented priors,
not a first-key rule that always dominates every other quality measure.

To claim the best candidate is optimal under $J_Q$, every unexplored search
state needs a valid lower bound on the same frozen objective, and the search
must not have discarded alternatives through incomplete pruning. The stopping
condition is then that the frontier lower bound is no smaller than the best
candidate cost, with a tie policy appropriate to the desired result set.
A grace-period cutoff yields only \emph{best seen}, not this certificate.

\imp Publish a total deterministic order: certification class, frozen cost,
and a canonical typed-expression tie-breaker. Evaluate held-out behavioral
accuracy, complexity, readability, and user preference separately. Reference
solution rank is useful, but allow proved equivalent implementations and more
than one reasonable intended solution.

\subsection{R8.2: which candidate equivalence is cheap and meaningful?}\label{q:R8.2}
\ans Use several explicitly labeled relations. No single cheap test provides
both general extensional equality and dependable canonicalization for all
dependent Lean terms.

The first layer is typed alpha-canonical structural identity, including the
relevant universes, binders, and context abstraction. The second is native
checked definitional equality. A third layer uses explicit equality proofs,
possibly pointwise proofs and function extensionality when permitted by the
profile. A fourth layer groups candidates that agree on a finite probe set.
Only the last layer should be called an \emph{observational cluster}, not a
proved equivalence class of programs.

The current engine's printed-form deduplication is useful for avoiding visual
duplicates, but the printer is not an identity function: it can hide implicit
arguments, universes, and coercions~\cite{implementation}. Replace its use in
semantic bookkeeping with an exact typed key. Keep readable-print equality
only as a display convenience, and make hidden distinctions available in the
artifact.

\begin{example}[Observational coincidence is not a general congruence]
Let $f(n)=n$ and let $g(n)=n$ for $n<2$, with $g(n)=0$ otherwise. They agree on
probes $\{0,1\}$. A surrounding program that applies its function argument to
$2$ distinguishes them. Therefore a synthesizer cannot replace $g$ by $f$ or
vice versa in arbitrary higher-order contexts merely because their current
probe signatures match.
\end{example}

The same issue arises in recursion when a surrounding program creates new
recursive inputs, or after CEGIS introduces a new counterexample. For a
restricted first-order straight-line DSL, observational quotienting can be
valid when every allowed operation acts pointwise on the retained observation
vectors and the relation is proved to be a congruence. The optimization is
valuable in that setting; its preconditions must not be omitted when moving
it into Lean's higher-order language.

Dependent functions introduce a further issue: observations at different
arguments may inhabit different fibers. Preserve the argument and type with
every observation, and compare outputs only by a well-typed relation. Proof
irrelevance can simplify proof components of the same proposition; it does not
remove the need to align data indices and transports.

\imp Keep one displayed representative per observational cluster while
retaining alternative implementations internally unless safe quotienting is
established. Search for a small distinguishing input between different
clusters and certify its output difference. Failure to find one within a bound
means only that the bounded test did not distinguish them, not that the
functions are equal.

\section{R9: learned proposals without expanding the trust boundary}\label{sec:r9}

\subsection{R9.1: what accuracy can be claimed for model-proposed carriers?}\label{q:R9.1}
\ans No numerical accuracy for the proposal's exact task is established by the
sources investigated here. Proof-retrieval results and general code-generation
benchmarks do not measure carrier or helper invention from only a Lean type
and finite examples. The appropriate answer is a controlled experiment with
precisely defined success criteria, not an extrapolated percentage.

Use the E8-style suite and the thirteen stretch cases as a starting point,
but avoid measuring only recovery of familiar textbook names. Include renamed
symbols, held-out function families, different encodings, and unfamiliar
combinations of known operations. Separate prompts containing only types and
examples from prompts containing a full symbolic specification or a reference
algorithm. The latter provide substantially different information.

For each model and fixed proposal budget, record four outcomes: well-formed
and well-typed proposal; a compatible implementation found under that proposal;
a kernel-certified solution to the original query; and improved end-to-end
quality or latency after charging model generation and validation. A proposed
carrier can be mathematically suitable yet difficult for the downstream
search, so carrier validity and end-to-end success should not be conflated.

Nor is exact agreement with a named reference carrier the right correctness
criterion. R2 shows that several representations can implement the same
function. Evaluate the complete carrier--seed--step--finisher solution under
the frozen objective, and report whether a carrier was merely plausible or
actually witnessed by a certified implementation.

Compare at least a deterministic grammar baseline, failure-driven symbolic
proposals, and the same search with model proposals added at fixed checkpoints.
Use identical downstream budgets and provider availability. Report per-family
outcomes and uncertainty rather than a single small-sample percentage. This
report does not run that model experiment, and therefore makes no claim about
current carrier-proposal accuracy.

\subsection{R9.2: what makes learned guidance deterministic and auditable?}\label{q:R9.2}
\ans Deterministic validation and replay of a frozen proposal artifact are more
important than hoping that model generation itself is bitwise deterministic.
A seed and temperature zero are not a sufficient reproducibility contract.

Record a content-addressed manifest containing the exact prompt input, model
and tokenizer identifiers, available runtime details, decoding parameters,
returned proposals, ordered alternatives, validation results, and the frozen
Lean environment and grammar fingerprints. Replay uses the recorded proposal
bytes and their parsed typed representations; it does not call the model
again. Canonical ordering and validation occur locally under the pinned
implementation.

The proposal language should be restricted data: carrier templates, references
to permitted constants, bounded helper signatures, or typed expression
skeletons. Do not treat arbitrary model-produced Lean commands or metaprograms
as safe merely because a final theorem will be checked. Elaboration-time code
can perform effects outside the mathematical kernel. A whitelist or a narrow
parser keeps proposal evaluation inside the intended boundary.

A model that only reorders an already fixed finite grammar preserves its
candidate universe, but it can still change budget-limited results. In an
infinite enumeration it can also cause starvation unless fairness is enforced.
If a model introduces a helper or carrier outside the recorded bounded
grammar, then the grammar has changed. Report that extension explicitly rather
than claiming unchanged bounded completeness. Alternatively, require every
accepted proposal to compile back into the original bounded grammar.

The audit record should connect each published term to its exact query,
contract proof, helper closure, axiom inventory, proposal manifest, and search
checkpoint. The mathematical proof remains valid independently of the model;
the receipt explains how the program was found and what search or ranking
claims are justified.

\imp Keep the default engine offline, core-only, and deterministic. An optional
proposal source should enter the same validated interface as enumerated
proposals and receive a bounded share of work. Do not add special acceptance
privileges for model-generated terms, and do not require the model to replay
an already certified result.

\section{An integrated Lean-native architecture}\label{sec:architecture}

\subsection{Keep authority at the original query}
The preceding answers do not require a second trusted type theory. Native
Lean expressions, local contexts, and metavariable states remain authoritative.
A separate residual or abstract representation is permitted for speed, but
its negative conclusions need a justified connection back to the native query.

The request-binding layer should retain the frozen target and contract rather
than accepting a caller-supplied pair of arbitrary types as sufficient context.
The inspected engine already constructs contract proofs by extracting the
property of a synthesized subtype; the gate then checks the supplied program
and proof types~\cite{implementation}. Preserve that connection explicitly in
the proposed interfaces, especially as more proposal sources and realization
paths are introduced.

There is a specific subtlety in the existing engine: its classical lane may
instantiate universe parameters at \code{Prop} and pass that specialized target
to acceptance~\cite{implementation}. An inhabitant of a specialization must be
labeled as such; it is not automatically an inhabitant of the original
universe-polymorphic request. This report does not claim that a kernel check
of the specialized term is invalid. It recommends making the distinction
visible in the result algebra and receipt.

\subsection{Proposed interfaces}
The following is an architectural schema, not compiled Lean source:
\begin{lstlisting}
FrozenQuery:
  target, contract, universePolicy, axiomProfile
  environmentFingerprint, grammarFingerprint, publicationObjective

Branch:
  savedLeanState, goals, dependentHoleTelescope
  residualGraph, pendingProofTasks, boundaryConstraints
  logicalCost, stableBranchId, proposalProvenance

Proposal:
  decisionKind, boundedSyntax, parentGrammar
  sourceManifest, validationEvidence

ObservationResult:
  proved(proof, dependencyReadSet)
  refuted(branchCertificate, dependencyReadSet)
  unknown(residual, blockers, reason)

PublishedResult:
  checkedProgram, checkedContractProof, exactRequestedType
  specializationLabel, axiomClosure, helperClosure
  frozenRank, searchCheckpoint, optimalityStatus
\end{lstlisting}

A \emph{closed candidate refutation} and an \emph{open branch refutation} should
be separate constructors. The first may be a proved counterexample to one
candidate. The second needs the completion-quantified statement in
\eqref{eq:refutation}. This prevents a common accidental promotion of a failed
guess into a claim about every future refinement.

Similarly, a proposal is not an accepted search state. It must first satisfy
scope, type, effect, grammar, and budget restrictions. A carrier proposal
creates ordinary seed, step, and finisher holes; a model-proposed helper must
be synthesized or supplied with a checked body, not added as an unexplained
constant.

\subsection{A deterministic control loop}
\begin{lstlisting}
freeze query, environment, grammar, objective, and proposal manifest
initialize a persistent frontier and a non-rollback work ledger
repeat until a logical bound is complete or cancellation occurs:
  choose the next deterministic scheduler checkpoint
  select a branch or proof continuation by recorded priority
  restore its Lean state
  validate caches and wake affected residual observations
  run cheap normalization and certified refutation
  if refuted: record evidence and discard this branch
  else if a complete candidate is available:
    realize recursive bodies and helpers, if needed
    check the original target, contract, and axiom profile
    update typed candidate identities and publication ranking
  else:
    compute conservative dependency components
    generate bounded native refinements and validated proposals
    add their saved successor states to the persistent frontier
  record completed work; never refund it on rollback
publish the best certified results and the exact completion status
\end{lstlisting}

This loop has no claim of universal completeness. Its bounded completeness is
parameterized by the recorded grammar, the fairness of refinement enumeration,
and the validity and completeness properties of the validators actually used.
A heuristic front end can run within it without corrupting the meaning of a
certified result.

\subsection{Negative-result algebra}
Four different negative-looking outcomes deserve distinct messages.
\emph{Contract impossible} requires a proof that no term of the requested type
satisfies the contract. \emph{Type uninhabited} requires a proof of its
negation under the recorded context and profile. \emph{Bounded grammar
exhausted} requires complete resolution of the stated finite search universe.
\emph{No result before cancellation} requires no mathematical conclusion at
all and must not be phrased as impossibility.

An unproved candidate that passes examples is another useful outcome, but it
belongs in a separate exploratory channel. It must not be mixed into the
kernel-certified candidate list with an ambiguous ``success'' label. The
same separation applies to a program that is mathematically certified but
whose executable binding has not yet been verified.

\section{Implementation sequence and falsifiable evaluation}\label{sec:roadmap}

\subsection{A dependency-aware sequence}
Do not wait for a general carrier-invention solution before improving all other
parts of the engine. Several high-value changes are independently testable.

\begin{longtable}{@{}L{12mm}L{62mm}L{77mm}@{}}
\toprule
Stage & Deliverable & Acceptance criterion\\
\midrule\endhead
P0 & Frozen request identity, typed candidate keys, exact outcome statuses & Existing accepted results recheck; specialization and cancellation are never mislabeled\\
P1 & Per-observation dependencies and rollback-safe caches & Cached and uncached decisions agree on the same saved states; sibling-branch and delayed-assignment tests pass\\
P2 & Persistent search frontier, deterministic scheduler, conservative components & Same logical checkpoint replays identically; no solved baseline case is lost at sufficient matched work\\
P3 & Native motive adapter, local invariant proofs, equation-based evaluation & Indexed map and supplied-carrier cases work without a replacement elaborator\\
P4 & Bounded joint carrier synthesis and finite-state refutations & Known witnesses are preserved; bounded exhaustion is reported only after all candidates are resolved\\
P5 & Capability bodies and progressively stronger angelic constraints & Every published recursive program has checked realization and contract evidence; fuel exhaustion remains unknown\\
P6 & Dictionary-aware retrieval and optional learned proposals & Provider leakage is controlled; proposal and retrieval overhead is charged; replay requires no model or network\\
\bottomrule
\end{longtable}

These are dependency stages, not promised delivery dates. P1 can begin with
the current evaluator; P3 can begin with a supplied carrier. The optional
learned layer should arrive after the deterministic proposal interface exists,
not become the reason that basic carrier and motive testing is deferred.

\subsection{The regression matrix}
A useful test suite should combine ordinary examples with deliberately
adversarial semantic cases.

\paragraph{Search state and caching.}
Include goals sharing a type metavariable, an index metavariable, a universe
constraint, and a dependency hidden in a local declaration. Include a tactic
that assigns a variable outside the apparent first goal. Test sibling branches
that reuse the same local version counter with different assignments. Compare
persistent-frontier and depth-first execution after sufficient exhaustive
work, not merely after the same wall-clock interval.

\paragraph{Carriers and recursors.}
Include reverse both with append and with a continuation, optional maximum,
min/max pairs, left folds with a non-associative operator, Fibonacci by tupling,
indexed vector map, an accumulating dependent function, and tree flattening
with and without append as a provider. Every benchmark should state the
permitted grammar and whether its reference witness lies within that grammar.
An absent provider is not evidence that a carrier is impossible.

\paragraph{Evaluation and proof.}
Test correct and deliberately unlawful \code{BEq} instances, partial terms
whose guards depend on holes, repeated recursive oracle calls, contradictory
observations, well-founded definitions requiring equations, and dependent
positions unsupported by a fast evaluator. For universal properties, include
one family whose invariant is supplied and another where it must be invented.
This separates proof-template execution from invariant discovery.

\paragraph{Ranking and retrieval.}
Include two expressions with the same pretty-printed display but different
hidden parameters, two candidates agreeing on current probes but differing
on an additional argument, legitimate constant functions, and components that
are useful only as intermediates. Include local typeclass instances that
change between contexts. Ensure target aliases and leaked reference bodies
cannot masquerade as synthesis achievements.

\subsection{Metrics and ablations}
Report representation coverage, first certified result, best frozen-objective
result, bounded-exhaustion rate, and unresolved/cancelled tasks separately.
Measure wall time as a practical outcome, but compare algorithmic policies
also at matched logical work. Track peak frontier memory: the proposal's cheap
snapshot time does not establish cheap retention of many live states.

For evaluation, count residual nodes reused, invalidated, and recomputed;
proof certificate construction and checking; and unsupported-fragment fallbacks.
For proof debt, count productive wakes and repeated attempts with unchanged
inputs. For retrieval, report candidate providers admitted, dictionary-solving
work, and the fraction of concrete type failures attributable to each abstract
feature. For carrier synthesis, record both the proposed carrier rank and the
rank of the first complete certified implementation under it.

Essential ablations are cache off/on; depth-first versus persistent frontier;
zero versus proved versus learned heuristic; carrier supplied versus invented;
first-order versus higher-order observations; native evaluation versus
proof-producing equations; and model off/on under the same total budget.
Changes that improve an aggregate solved count while creating incorrect hard
pruning should fail acceptance immediately.

\subsection{What remains genuinely empirical}
This analysis does not settle the best weights of a carrier grammar, the
most useful abstraction depth for Mathlib retrieval, the best proof-debt
scheduler, or current model accuracy on the exact requested task. Those
quantities require experiments on a pinned corpus and implementation. The
report instead provides well-defined experiments and interfaces that make
results interpretable. In particular, it replaces unqualified claims of
minimality, completeness, and reproducibility with statements whose premises
can be checked.

\section{Conclusion}
Most of the expert questions do not require abandoning Leant2's Lean-native
foundation. They require sharper boundaries between a proposal, an observation,
a refutation, a proof, and a search-order preference. The strongest practical
route is a persistent native search with conservative dependency factoring,
rollback-safe observations, checked native motive construction, local
invariant obligations, and a bounded joint carrier language. Existing work
already supplies higher-order live examples, dictionary-aware synthesis, and
proof-producing evaluation techniques that materially shorten the research
path.

The hardest remaining problem is not naming a clever carrier in isolation.
It is finding an economical, type-correct combination of state, transitions,
finishing computation, recursive scheme, and invariant under a realistic
budget. Treating that combination as a bounded, evidence-producing synthesis
problem gives useful positive and negative answers without overstating what
finite examples, a kernel gate, or learned guidance can establish.

\appendix
\section{Complete question-to-answer index}\label{app:map}
\begingroup
\renewcommand{\arraystretch}{1.02}
\small
\begin{longtable}{@{}L{14mm}L{126mm}r@{}}
\toprule
ID & Question focus & Page\\
\midrule\endhead
R1.1 & Shared metavariables, independence, and cross-branch memoization & \pageref{q:R1.1}\\
R1.2 & Admissible cost-to-close heuristic & \pageref{q:R1.2}\\
R1.3 & Just-in-time learning and ranking & \pageref{q:R1.3}\\
R1.4 & Machine-speed independence under wall-clock budgets & \pageref{q:R1.4}\\
R2.1 & Minimal grammar carrier from finite observations & \pageref{q:R2.1}\\
R2.2 & Polymorphic alphabet and parametricity & \pageref{q:R2.2}\\
R2.3 & Backward fusion, automata, and useful complete classes & \pageref{q:R2.3}\\
R2.4 & Motive enumeration for parameterized primitive recursion & \pageref{q:R2.4}\\
R2.5 & Generalization, rippling, and lemma speculation & \pageref{q:R2.5}\\
R3.1 & Higher-order and dependent live bidirectional evaluation & \pageref{q:R3.1}\\
R3.2 & Assignment-driven normalization and rollback & \pageref{q:R3.2}\\
R3.3 & A fast conservative evaluator and Lean4Lean & \pageref{q:R3.3}\\
R3.4 & Top-down angelic termination and completeness & \pageref{q:R3.4}\\
R4.1 & Complete bounded motive classes and accumulators & \pageref{q:R4.1}\\
R4.2 & Canonical treatment of transports, setoids, and heterogeneous equality & \pageref{q:R4.2}\\
R4.3 & Definitional unfolding policy & \pageref{q:R4.3}\\
R4.4 & Lean motive construction as a library & \pageref{q:R4.4}\\
R5.1 & Capability interpreter versus realized recursion & \pageref{q:R5.1}\\
R5.2 & Angelic execution in top-down type-directed search & \pageref{q:R5.2}\\
R5.3 & A small recursion basis and measured coverage & \pageref{q:R5.3}\\
R6.1 & Interleaving value and proof holes without a universal SMT encoding & \pageref{q:R6.1}\\
R6.2 & A recognizable sufficient class of inductive proofs & \pageref{q:R6.2}\\
R6.3 & Current Lean induction automation and bounded use & \pageref{q:R6.3}\\
R6.4 & Principled scheduling of proof debt & \pageref{q:R6.4}\\
R7.1 & Dependent and universe-polymorphic reachability abstraction & \pageref{q:R7.1}\\
R7.2 & Type classes in type-guided abstraction refinement & \pageref{q:R7.2}\\
R7.3 & Data-goal premise selection and training corpora & \pageref{q:R7.3}\\
R8.1 & A formal objective for ranking and its empirical validation & \pageref{q:R8.1}\\
R8.2 & Cheap meaningful equivalence for dependent candidates & \pageref{q:R8.2}\\
R9.1 & Accuracy of learned carrier/helper proposals & \pageref{q:R9.1}\\
R9.2 & Deterministic and auditable model contribution & \pageref{q:R9.2}\\
\bottomrule
\end{longtable}
\endgroup

\section{Executable semantic checks and their limits}\label{app:checks}
The script \code{semantic_checks.py} uses only Python's standard library.
Its recorded output is \code{semantic_results.json}.
These checks are small semantic experiments, not execution of Lean or Leant2.
They neither establish performance improvements nor replace the mathematical
arguments in the report.

\begin{longtable}{@{}L{49mm}L{66mm}L{36mm}@{}}
\toprule
Check & Scope & Observed result\\
\midrule\endhead
Reverse via ordinary fold and continuation & All 3,280 lists over $\{0,1,2\}$ of lengths zero through seven & Both match reversal\\
Optional min/max state & The same 3,280 lists, including empty input & All match reference\\
Fibonacci pair update & Indices zero through sixty; reference uses fast doubling & All 61 match\\
Unary finite-state synthesis & Boolean outputs for lengths zero through eight; target length divisible by three & Zero consistent 1- or 2-state machines; two labeled 3-state machines\\
Residual-state lower bound & Three tails and three prefix continuations for the unary target & An incompatibility clique of size three\\
Incomplete observation rows & Two different partial rows with no conflicting common column & Rows are compatible\\
Observational congruence & Two functions equal on probes zero and one & Application at two distinguishes them\\
Rollback cache identity & Two sibling assignments with reused version number & Naive key fails; qualified key separates them\\
Goal-count lower bound & Abstract two-obligation shared-assignment example & Remaining cost one is below goal count two\\
\bottomrule
\end{longtable}

For the finite-state experiment, the initial state is fixed to zero. With $k$
states, a unary transition table has $k^k$ possibilities and a Boolean output
table has $2^k$. The exhaustive enumeration therefore checks $2$, $16$, and
$216$ machines for $k=1,2,3$. The two successful three-state tables differ by
state labeling. The clique witnesses are tail lengths $0,1,2$, tested after
prefix lengths $0,1,2$; their fully known output rows are pairwise incompatible.
This supplies a genuine finite-state lower bound rather than a count of
syntactically distinct partial rows.

The reverse checks use Python list operations and are \emph{not} evidence for
Lean runtime complexity. The asymptotic discussion in R2 is about the stated
linked-list model. The cache and heuristic examples exercise the logical
countermodels; they are not demonstrations of a confirmed bug in the current
Leant2 implementation.

\clearpage
\section{Evidence status and reproducibility}\label{app:status}
The primary research sources and official Lean documentation were consulted
on 21 September 2026. Repository analysis is pinned to the SHA on the title
page; native API observations are pinned to Lean 4.34. The proposal's older
benchmark measurements are cited as historical reports. The mathematical
arguments in this report are written proofs, not Lean-formalized theorems.

No Lean executable was available in the working environment, so none of the
proposed Lean interfaces or examples was compiled here. The report does not
claim a Leant2 patch, a rerun of its acceptance harnesses, cross-machine replay
measurements, or model accuracy data. The Python semantic checks were executed
successfully, and the accompanying PDF was compiled from this source.

The inspected material did not establish a precise current \code{sauto}
unfolding policy or a published accuracy measurement matching R9.1's exact
carrier-proposal task. The relevant sections give explicit recommendations
and evaluation procedures without filling those evidence gaps with invented
results. Claims of absence from the literature are deliberately avoided;
``not established by the sources investigated'' is narrower and more accurate.

To rebuild, run \code{pdflatex leant2_expert_answers.tex} twice, or run
\code{latexmk -pdf leant2_expert_answers.tex}. A normal TeX Live installation
with the packages named in the preamble is sufficient. Run
\code{python semantic_checks.py} to reproduce the semantic result file.

\clearpage
\begin{thebibliography}{99}
\small

\bibitem{repo}
V.~Reshetnikov. \emph{Leant2: Further Improvements}, question-bearing sections
05--12 and expert summary 15. Repository snapshot
\code{784664ff34e819422510a4d470ab07fe48bb736b}, inspected 21 September 2026.
\url{https://github.com/VladimirReshetnikov/Leant2/tree/784664ff34e819422510a4d470ab07fe48bb736b/docs/proposals/11-further-improvements}.

\bibitem{implementation}
V.~Reshetnikov. \emph{Leant2 implementation}: \code{Leant2/Engine.lean},
\code{Leant2/Native/Transaction.lean}, \code{Leant2/Accept/Gate.lean}, and
\code{lean-toolchain}. Same pinned snapshot as~\cite{repo}.
\href{https://github.com/VladimirReshetnikov/Leant2/blob/784664ff34e819422510a4d470ab07fe48bb736b/Leant2/Engine.lean}{Engine source};
\href{https://github.com/VladimirReshetnikov/Leant2/blob/784664ff34e819422510a4d470ab07fe48bb736b/Leant2/Native/Transaction.lean}{transaction source};
\href{https://github.com/VladimirReshetnikov/Leant2/blob/784664ff34e819422510a4d470ab07fe48bb736b/Leant2/Accept/Gate.lean}{gate source}.

\bibitem{proposal-evidence}
V.~Reshetnikov. \emph{Leant2 Further Improvements: Evidence},
\code{sections/01-evidence.tex}. Same pinned snapshot as~\cite{repo}; the
acceptance table itself reports commit \code{33cec8b}.
\href{https://github.com/VladimirReshetnikov/Leant2/blob/784664ff34e819422510a4d470ab07fe48bb736b/docs/proposals/11-further-improvements/sections/01-evidence.tex}{Pinned evidence file}.

\bibitem{aesop}
J.~Limperg and A.~H.~From. An Extensible Theorem Prover for Lean.
\emph{CPP}, 2023. DOI: \href{https://doi.org/10.1145/3573105.3575671}{10.1145/3573105.3575671}.
\url{https://people.compute.dtu.dk/ahfrom/aesop-camera-ready.pdf}.

\bibitem{canonical}
C.~Norman and J.~Avigad. Canonical for Automated Theorem Proving in Lean.
2025. \url{https://arxiv.org/abs/2504.06239}.

\bibitem{canonical-min}
C.~Norman and J.~Avigad. Implementing Dependent Type Theory Inhabitation and
Unification. 2026. \url{https://arxiv.org/abs/2603.01463}.

\bibitem{probe}
S.~Barke, H.~Peleg, and N.~Polikarpova. Just-in-Time Learning for Bottom-Up
Enumerative Synthesis. \emph{PACMPL}, OOPSLA, 2020.
DOI: \href{https://doi.org/10.1145/3428295}{10.1145/3428295}.
\url{https://arxiv.org/abs/2010.08663}.

\bibitem{fold-characterization}
J.~Gibbons, G.~Hutton, and T.~Altenkirch. When Is a Function a Fold or an
Unfold? \emph{Electronic Notes in Theoretical Computer Science}, 44(1), 2001.
\url{https://www.cs.ox.ac.uk/publications/publication2368-abstract.html}.

\bibitem{constructive-fold}
T.~Weber and J.~Caldwell. Constructively Characterizing Fold and Unfold.
\emph{LOPSTR 2003, Revised Selected Papers}, LNCS 3018, 2004.
DOI: \href{https://doi.org/10.1007/978-3-540-25938-1_11}{10.1007/978-3-540-25938-1\_11}.
\url{https://www.cs.uwyo.edu/~jlc/papers/fold_unfold.pdf}.

\bibitem{angluin}
D.~Angluin. Learning Regular Sets from Queries and Counterexamples.
\emph{Information and Computation}, 75:87--106, 1987.
DOI: \href{https://doi.org/10.1016/0890-5401(87)90052-6}{10.1016/0890-5401(87)90052-6}.

\bibitem{parametric-testing}
J.-P.~Bernardy, P.~Jansson, and K.~Claessen. Testing Polymorphic Properties.
\emph{ESOP}, LNCS 6012, pp.~125--144, 2010.
DOI: \href{https://doi.org/10.1007/978-3-642-11957-6_8}{10.1007/978-3-642-11957-6\_8}.
\url{https://publications.lib.chalmers.se/records/fulltext/local_99387.pdf}.

\bibitem{hipspec}
K.~Claessen, M.~Johansson, D.~Ros\'en, and N.~Smallbone. Automating Inductive
Proofs Using Theory Exploration. \emph{CADE-24}, LNCS 7898, pp.~392--406, 2013.
DOI: \href{https://doi.org/10.1007/978-3-642-38574-2_27}{10.1007/978-3-642-38574-2\_27}.
\url{https://www.cse.chalmers.se/~jomoa/papers/hipspec-cade.pdf}.

\bibitem{smyth}
J.~Lubin, N.~Collins, C.~Omar, and R.~Chugh. Program Sketching with Live
Bidirectional Evaluation. \emph{PACMPL}, 4(ICFP), article 109, 2020.
DOI: \href{https://doi.org/10.1145/3408991}{10.1145/3408991}.
Extended version: \url{https://arxiv.org/abs/1911.00583}.

\bibitem{adapton}
M.~A.~Hammer, K.~Y.~Phang, M.~Hicks, and J.~S.~Foster. Adapton: Composable,
Demand-Driven Incremental Computation. \emph{PLDI}, 2014.
DOI: \href{https://doi.org/10.1145/2594291.2594324}{10.1145/2594291.2594324}.
\url{https://www.cs.umd.edu/~mwh/papers/hammer13adapton.html}.

\bibitem{lean-cbv}
Lean FRO and contributors. \emph{The Lean Language Reference}, version 4.34.0,
Tactic Reference: \code{cbv}, \code{decide_cbv}, and their limitations.
\url{https://lean-lang.org/doc/reference/4.34.0/Tactic-Proofs/Tactic-Reference/}.

\bibitem{lean-recursion}
Lean FRO and contributors. \emph{The Lean Language Reference}: Recursive
Definitions. Online reference consulted 21 September 2026.
\url{https://lean-lang.org/doc/reference/latest/Definitions/Recursive-Definitions/}.

\bibitem{lean4lean}
M.~Carneiro. Lean4Lean: Towards a Verified Typechecker for Lean, in Lean.
\emph{TYPES 2025}, LIPIcs 384, article 2, published 30 July 2026.
DOI: \href{https://doi.org/10.4230/LIPIcs.TYPES.2025.2}{10.4230/LIPIcs.TYPES.2025.2}.
\url{https://drops.dagstuhl.de/entities/document/10.4230/LIPIcs.TYPES.2025.2}.

\bibitem{burst}
A.~Miltner, A.~T.~Nu\~nez, A.~Brendel, S.~Chaudhuri, and I.~Dillig.
Bottom-Up Synthesis of Recursive Functional Programs Using Angelic Execution.
\emph{PACMPL}, 6(POPL), article 21, 2022.
DOI: \href{https://doi.org/10.1145/3498682}{10.1145/3498682}.
\url{https://arxiv.org/abs/2107.06253}.

\bibitem{sauto}
\L.~Czajka. Practical Proof Search for Coq by Type Inhabitation.
\emph{IJCAR}, pp.~28--57, 2020.
DOI: \href{https://doi.org/10.1007/978-3-030-51054-1_3}{10.1007/978-3-030-51054-1\_3}.
Current implementation-level unfolding policy was not verified in this report.

\bibitem{lean-induction-source}
Lean contributors. \code{src/Lean/Meta/Tactic/Induction.lean}, version
\code{v4.34.0}; native induction, index extraction, and recursor-prefix APIs.
\url{https://github.com/leanprover/lean4/blob/v4.34.0/src/Lean/Meta/Tactic/Induction.lean}.

\bibitem{synquid}
N.~Polikarpova, I.~Kuraj, and A.~Solar-Lezama. Program Synthesis from
Polymorphic Refinement Types. \emph{PLDI}, 2016.
DOI: \href{https://doi.org/10.1145/2908080.2908093}{10.1145/2908080.2908093}.
\url{https://cseweb.ucsd.edu/~npolikarpova/publications/pldi16.pdf}.

\bibitem{leanhammer}
T.~Zhu, J.~Clune, J.~Avigad, A.~Q.~Jiang, and S.~Welleck. Premise Selection
for a Lean Hammer. \emph{ICLR}, 2026; preprint first posted in 2025.
\url{https://arxiv.org/abs/2506.07477}.

\bibitem{hoogle}
Z.~Guo, M.~James, D.~Justo, J.~Zhou, Z.~Wang, R.~Jhala, and N.~Polikarpova.
Program Synthesis by Type-Guided Abstraction Refinement.
\emph{PACMPL}, 4(POPL), article 12, 2020.
DOI: \href{https://doi.org/10.1145/3371080}{10.1145/3371080}.
\url{https://ranjitjhala.github.io/static/hoogle_plus.pdf}.

\bibitem{lean-selector-source}
Lean contributors. \code{src/Lean/LibrarySuggestions/Basic.lean}, version
\code{v4.34.0}; selector and configuration API.
\url{https://github.com/leanprover/lean4/blob/v4.34.0/src/Lean/LibrarySuggestions/Basic.lean}.

\bibitem{selector-readme}
Premise-selection project contributors. \emph{premise-selection README},
consulted 21 September 2026; cloud selector interface, caching, and cold-start
behavior. Inspected README blob:
\code{ba823550133d3b56a756787eab334f8151e03808}.
\url{https://github.com/hanwenzhu/premise-selection/blob/main/README.md}.

\bibitem{ranking-study}
N.~Natarajan, D.~Simmons, N.~Datha, P.~Jain, and S.~Gulwani.
Learning Natural Programs from a Few Examples in Real-Time.
\emph{AISTATS}, PMLR 89:1714--1722, 2019.
\url{https://proceedings.mlr.press/v89/natarajan19a.html}.

\bibitem{pbe-interaction}
M.~Mayer et al. User Interaction Models for Disambiguation in Programming
by Example. \emph{UIST}, 2015.
DOI: \href{https://doi.org/10.1145/2807442.2807459}{10.1145/2807442.2807459}.
\url{https://www.microsoft.com/en-us/research/publication/user-interaction-models-disambiguation-programming-example/}.

\end{thebibliography}
\end{document}
